% ============================================================
%  Příprava na maturitu – Mechatronika
%  Kompiluj: pdflatex maturita_mechatronika.tex
%  Kódování: UTF-8
% ============================================================

\documentclass[a4paper,10pt]{article}

% --- Balíčky ---
\usepackage[T1]{fontenc}
\usepackage[utf8]{inputenc}
\usepackage[czech]{babel}
\usepackage[left=2.4cm, right=1.8cm, top=1.8cm, bottom=1.8cm]{geometry}
\usepackage{lmodern}
\usepackage{microtype}
\usepackage{parskip}
\usepackage{titlesec}
\usepackage{enumitem}
\usepackage{xcolor}
\usepackage{hyperref}
\usepackage{tabularx}
\usepackage{booktabs}
\usepackage{array}
\usepackage{tikz}
\usepackage{eso-pic}
\usepackage{mdframed}
\usepackage{amsmath}
\usepackage{amssymb}
\usepackage{pgfplots}
\usepackage{karnaugh-map}
\usepackage{circuitikz}
\pgfplotsset{compat=1.18}
\usetikzlibrary{arrows.meta, positioning, shapes.geometric, calc}

% --- Barvy ---
\definecolor{accent}{HTML}{03254E}
\definecolor{stripeblue}{HTML}{568EA3}
\definecolor{medgray}{HTML}{888888}
\definecolor{lightblue}{HTML}{EAF2F8}
\definecolor{lightyellow}{HTML}{FDFBE4}
\definecolor{lightgreen}{HTML}{EAF7EA}

% --- Modrý pruh vlevo ---
\AddToShipoutPictureBG{%
  \begin{tikzpicture}[remember picture, overlay]
    \fill[stripeblue]
      (current page.north west)
      rectangle
      ([xshift=10mm] current page.south west);
    \fill[accent!60]
      ([xshift=10mm] current page.north west)
      rectangle
      ([xshift=12mm] current page.south west);
  \end{tikzpicture}%
}

\hypersetup{colorlinks=true, urlcolor=accent, linkcolor=accent}

\titleformat{\section}
  {\large\bfseries\color{accent}}
  {\thesection.}
  {0.5em}
  {}
  [\color{accent}\titlerule]

\titlespacing{\section}{0pt}{16pt}{6pt}

\newmdenv[
  backgroundcolor=lightblue,
  linecolor=stripeblue,
  linewidth=1pt,
  roundcorner=4pt,
  innerleftmargin=8pt,
  innerrightmargin=8pt,
  innertopmargin=6pt,
  innerbottommargin=6pt,
  skipabove=4pt,
  skipbelow=4pt
]{pojmybox}

\newmdenv[
  backgroundcolor=lightyellow,
  linecolor=accent!40,
  linewidth=1pt,
  roundcorner=4pt,
  innerleftmargin=8pt,
  innerrightmargin=8pt,
  innertopmargin=6pt,
  innerbottommargin=6pt,
  skipabove=4pt,
  skipbelow=8pt
]{vysvetlenibox}

\newmdenv[
  backgroundcolor=lightgreen,
  linecolor=accent!30,
  linewidth=1pt,
  roundcorner=4pt,
  innerleftmargin=8pt,
  innerrightmargin=8pt,
  innertopmargin=6pt,
  innerbottommargin=6pt,
  skipabove=4pt,
  skipbelow=8pt
]{vzorcebox}

\pagestyle{plain}

% ============================================================
\begin{document}

% --- Záhlaví (pouze na první straně) ---
\begin{minipage}[t]{0.75\textwidth}
    {\Huge\bfseries\color{accent} Příprava na maturitu}\\[4pt]
    {\large\color{stripeblue} Mechatronika}\\[2pt]
\end{minipage}

% ============================================================
\section{Mechatronika, mechatronický návrh, systém a výrobek}

\begin{pojmybox}
\textbf{Klíčové pojmy:}
mechatronika, mechanika, elektronika, informatika,
senzor, aktuátor, řídící systém, systémový přístup, mechatronický výrobek,
V-model návrhu, SiL, HiL, zpětná vazba
\end{pojmybox}

\begin{vysvetlenibox}
\textbf{Definice a vznik oboru}\\
Mechatronika je interdisciplinární obor kombinující \textbf{mechaniku},
\textbf{elektroniku} a \textbf{informatiku} za účelem návrhu inteligentních
systémů a výrobků. Termín poprvé použila japonská firma Yaskawa Electric
v roce 1969. Dnes je mechatronika základem moderní průmyslové automatizace,
automobilového průmyslu, robotiky i spotřební elektroniky.

\textbf{Historický vývoj mechatroniky}\\
Vývoj mechatroniky lze rozdělit do několika fází:
\begin{itemize}[noitemsep, leftmargin=1.4em, topsep=4pt]
    \item \textbf{1969--1980} -- vznik termínu, první mechatronické systémy (přádací stroje, CNC)
    \item \textbf{1980--1995} -- rozvoj mikroprocesorů, masové nasazení v automobilkách
    \item \textbf{1995--2010} -- nástup průmyslových sběrnic, decentralizovaná řízení
    \item \textbf{2010--současnost} -- Industry 4.0, IoT, umělá inteligence, kyberfyzikální systémy
\end{itemize}

\textbf{Složky mechatronického systému}\\
Každý mechatronický výrobek obsahuje tři neoddělitelné složky:
\begin{itemize}[noitemsep, leftmargin=1.4em, topsep=4pt]
    \item \textbf{Mechanická část} -- nosná konstrukce, převody, kinematické
          vazby; zajišťuje pohyb a přenos sil
    \item \textbf{Elektronická část} -- senzory (měří fyzikální veličiny),
          aktuátory (vykonávají pohyb), výkonová elektronika (řídí výkon)
    \item \textbf{Řídicí část} -- mikrokontrolér nebo PLC s algoritmem;
          zpracovává data ze senzorů a vydává příkazy aktuátorům
\end{itemize}

\textbf{Informační toky v mechatronickém systému}\\
Signál putuje systémem v uzavřené smyčce:
\begin{enumerate}[noitemsep, leftmargin=1.4em, topsep=4pt]
    \item Fyzikální veličina (teplota, tlak, poloha) působí na \textbf{senzor}
    \item Senzor převede veličinu na elektrický signál (napětí, proud, frekvence)
    \item Signál je kondicionován (zesílen, filtrován) a převeden A/D převodníkem
    \item \textbf{Řídicí jednotka} vyhodnotí stav a vypočítá akční zásah
    \item Akční člen (\textbf{aktuátor}) provede pohyb nebo změnu výkonu
    \item Nový stav je opět měřen senzorem -- smyčka se uzavírá
\end{enumerate}

\textbf{Systémový přístup a V-model}\\
Na rozdíl od tradičního sekvenčního návrhu (nejprve mechanika, pak elektronika,
nakonec software) systémový přístup navrhuje všechny složky \textit{současně}.
Standardním postupem je \textbf{V-model}:

\vspace{4pt}
\begin{center}
\begin{tikzpicture}[scale=0.85,
  box/.style={rectangle, draw=accent, fill=lightblue, rounded corners=3pt,
              minimum width=2.8cm, minimum height=0.6cm, font=\small,
              text centered}]
  \node[box] (spec)    at (0,0)    {Specifikace požadavků};
  \node[box] (arch)    at (2,-1)   {Architektura systému};
  \node[box] (detail)  at (4,-2)   {Detailní návrh};
  \node[box] (impl)    at (6,-3)   {Implementace};
  \node[box] (unit)    at (8,-2)   {Unit testování};
  \node[box] (integ)   at (10,-1)  {Integrační testy};
  \node[box] (valid)   at (12,0)   {Validace systému};
  \draw[>-, accent, thick]
    (spec) -- (arch) -- (detail) -- (impl) -- (unit) -- (integ) -- (valid);
  \draw[dashed, medgray]
    (spec.south east) -- (valid.south west)
    node[midway, below, font=\scriptsize, color=medgray] {čas / fáze};
\end{tikzpicture}
\end{center}
\end{vysvetlenibox}

\begin{vysvetlenibox}
\textbf{Fáze V-modelu v praxi:}
\begin{itemize}[noitemsep, leftmargin=1.4em, topsep=4pt]
    \item \textbf{Specifikace požadavků} -- co má systém umět (rychlost, přesnost, bezpečnost)
    \item \textbf{Architektura systému} -- rozdělení na mechaniku, elektroniku, software
    \item \textbf{Detailní návrh} -- CAD modely, schémata elektroniky, algoritmy
    \item \textbf{Implementace} -- výroba prototypu, programování, osazování
    \item \textbf{Unit testování} -- testování jednotlivých komponent
    \item \textbf{Integrační testy} -- ověření spolupráce všech složek
    \item \textbf{Validace systému} -- splnění původních požadavků zákazníka
\end{itemize}

\textbf{SiL a HiL testování}\\
Moderní vývoj bez fyzického prototypu v raných fázích:
\begin{itemize}[noitemsep, leftmargin=1.4em, topsep=4pt]
    \item \textbf{SiL} (Software in the Loop) -- software řídicího algoritmu
          běží na PC a simuluje chování soustavy; testuje logiku bez HW.
          \textit{Příklad:} Simulace PID regulátoru v MATLAB/Simulink.
    \item \textbf{HiL} (Hardware in the Loop) -- reálný řídicí HW (PLC, MCU)
          je připojen k simulátoru fyzikální soustavy; testuje se odezva na
          reálná vstupní data v reálném čase.
          \textit{Příklad:} Ověření airbag řídicí jednotky s simulátorem nárazu.
    \item \textbf{MiL} (Model in the Loop) -- testování matematického modelu
          proti požadavkům; nejranější fáze verifikace.
\end{itemize}

\textbf{Příklady mechatronických výrobků}

\begin{center}
\begin{tabularx}{0.95\linewidth}{@{} l X X X @{}}
    \toprule
    \textbf{Výrobek} & \textbf{Mechanika} & \textbf{Elektronika}
                     & \textbf{Řídicí systém} \\
    \midrule
    ABS brzdy        & brzdový třmen, pístnice & senzory otáček kol
                     & ECU, PID regulátor \\
    Průmyslový robot & kinematický řetězec, klouby & servomotory, enkodéry
                     & trajektoriový kontrolér \\
    Pračka           & buben, ložiska, čerpadlo & motor, ventily, senzory
                     & PLC / MCU, program cyklů \\
    CNC frézka       & portál, vřeteno, osy & servopohony, čidla polohy
                     & CNC řídicí systém, G-kód \\
    \bottomrule
\end{tabularx}
\end{center}

\textbf{Detailní příklad: Antiblokovací brzdový systém (ABS)}\\
ABS je ukázkový mechatronický systém v automobilu:
\begin{itemize}[noitemsep, leftmargin=1.4em, topsep=4pt]
    \item \textbf{Mechanika:} Brzdový třmen tlačí brzdové destičky na kotouč.
          Pístnice přenáší sílu z brzdového pedálu.
    \item \textbf{Elektronika:} Indukční senzory u každého kola měří okamžitou
          úhlovou rychlost. Při rozdílu otáček kol detekují smyk.
    \item \textbf{Řízení:} Řídicí jednotka (ECU) vyhodnocuje signály ze senzorů
          100× za sekundu. Při detekci smyku otevře elektromagnetický ventil,
          který sníží tlak v brzdovém systému a kolo se opět roztočí.
\end{itemize}

\textbf{Výhody mechatronického přístupu:}
\begin{itemize}[noitemsep, leftmargin=1.4em, topsep=4pt]
    \item \textbf{Optimalizace celku} -- lze najít lepší řešení než optimalizací každé části zvlášť
    \item \textbf{Flexibilita} -- změny ve funkci se často řeší úpravou software
    \item \textbf{Diagnostika} -- systém umí detekovat poruchy a informovat obsluhu
    \item \textbf{Adaptivita} -- moderní systémy se učí a přizpůsobují podmínkám
\end{itemize}

\textbf{Nevýhody a výzvy:}
\begin{itemize}[noitemsep, leftmargin=1.4em, topsep=4pt]
    \item \textbf{Komplexita} -- těžší pochopení a ladění celého systému
    \item \textbf{Bezpečnost} -- software může obsahovat chyby s vážnými důsledky
    \item \textbf{Kyberbezpečnost} -- připojené systémy jsou zranitelné vůči útokům
    \item \textbf{Odborníci} -- potřeba inženýrů s přesahem do více oborů
\end{itemize}
\end{vysvetlenibox}

% ============================================================

% ============================================================

\newpage

% ============================================================
\section{Logické řízení -- programovací jazyky, logické funkce}

\begin{pojmybox}
\textbf{Klíčové pojmy:}
Booleova algebra, Karnaughova mapa, logické funkce (AND, OR, NOT, NAND, NOR, XOR),
pravdivostní tabulka, funkčně úplný systém
\end{pojmybox}

\begin{vysvetlenibox}
	\textbf{Úvod do logiky}
	\begin{itemize}[noitemsep]
		\item z řeckého slova logos $\Rightarrow$ důvod
		\item Věda zkoumající způsob vyvozování závěrů a způsob předpokladu
		\item Myšlenková cesta, která vede k daným závěrům
	\end{itemize}
	

	\textbf{Výrok}
	\begin{itemize}[noitemsep]
		\item Tvrzení u kterého má smysl ptát se, zda je \textbf{pravdivé} či \textbf{nepravdivé}
		\item Jednoduchý výrok $\rightarrow$ "Včera tady pršelo"
		\item Složený výrok $\rightarrow$ "Včera tady pršelo a současně se jmenuji Pepa"
	\end{itemize}
	
\textbf{Logické spojky}

\begin{center}
	\resizebox{\textwidth}{!}{
		\begin{tabular}{|l|l|p{4cm}|p{5.5cm}|}
			\hline
			\textbf{Operace} & \textbf{Zápis} & \textbf{Význam} & \textbf{Příklad} \\
			\hline
			Negace & $\neg A$, NOT A & není pravda A & $\neg(\text{Prší}) = \text{Neprší}$ \\
			\hline
			Konjunkce & $A \land B$, $A \cdot B$, A AND B & A a současně B & Nemám vlasy a současně se jmenuji Pepa \\
			\hline
			Disjunkce & $A \lor B$, A + B, A OR B & A nebo B & Vlak přijel na 1. nebo 2. kolej \\
			\hline
			Implikace & $A \Rightarrow B$ & jestliže A, pak B & Jestliže pršelo, je mokro \\
			\hline
			Ekvivalence & $A \Leftrightarrow B$ & A tehdy a jen tehdy B & Číslo je sudé tehdy a jen tehdy, pokud je dělitelné 2 \\
			\hline
		\end{tabular}
	}
\end{center}

\textbf{Důležité pojmy}
\begin{itemize}[noitemsep]
	\item funkce, definiční obor a obor hodnot
	\item logická proměnná
	\item hradlo
	\item logický obvod
\end{itemize}
	
\textbf{Booleova algebra}\\
Logické řízení pracuje s binárními signály: \textbf{0} (nepravda, vypnuto)
a \textbf{1} (pravda, zapnuto). Základem je Booleova algebra -- algebraická
struktura s operacemi součin (AND), součet (OR) a negace (NOT).

\textbf{Fyzikální reprezentace logických úrovní}
\begin{itemize}[noitemsep]
    \item \textbf{TTL logika} (Transistor-Transistor Logic):
          $0 = 0\text{--}0{,}8\,\text{V}$, $1 = 2\text{--}5\,\text{V}$
    \item \textbf{CMOS logika} (Complementary MOS):
          $0 = 0\text{--}1{,}5\,\text{V}$, $1 = 3{,}5\text{--}5\,\text{V}$
    \item \textbf{PLC vstupy} (podle normy IEC 61131-2):
          $0 = 0\text{--}5\,\text{V}$, $1 = 15\text{--}30\,\text{V}$
\end{itemize}

\textbf{Základní logické funkce a jejich pravdivostní tabulky}

\begin{center}
\begin{tabular}{cc|c|c|c|c|c}
    \toprule
    $A$ & $B$ & AND & OR & NAND & NOR & XOR \\
    \midrule
    0 & 0 & 0 & 0 & 1 & 1 & 0 \\
    0 & 1 & 0 & 1 & 1 & 0 & 1 \\
    1 & 0 & 0 & 1 & 1 & 0 & 1 \\
    1 & 1 & 1 & 1 & 0 & 0 & 0 \\
    \bottomrule
\end{tabular}
\end{center}

\textbf{Vysvětlení logických funkcí:}
\begin{itemize}[noitemsep, leftmargin=1.4em, topsep=4pt]
    \item \textbf{AND} (konjunkce, $\cdot$) -- výstup je 1, pouze když \textit{oba} vstupy jsou 1
    \item \textbf{OR} (disjunkce, $+$) -- výstup je 1, když \textit{alespoň jeden} vstup je 1
    \item \textbf{NOT} (negace, $\overline{A}$) -- převrátí hodnotu (0$\to$1, 1$\to$0)
    \item \textbf{NAND} (NOT-AND) -- negovaný AND; univerzální hradlo
    \item \textbf{NOR} (NOT-OR) -- negovaný OR; také univerzální
    \item \textbf{XOR} (exklusivní OR) -- výstup je 1, když se vstupy \textit{liší}
\end{itemize}

\textbf{Základní zákony Booleovy algebry}

\begin{vzorcebox}
\begin{align*}
A \cdot 0 &= 0 & A + 0 &= A & \text{(nulový prvek)} \\
A \cdot 1 &= A & A + 1 &= 1 & \text{(jednotkový prvek)} \\
A \cdot A &= A & A + A &= A & \text{(idempotence)} \\
A \cdot \overline{A} &= 0 & A + \overline{A} &= 1 & \text{(komplementarita)} \\
\overline{\overline{A}} &= A & & & \text{(dvojí negace)} \\
\overline{A \cdot B} &= \overline{A} + \overline{B} &
\overline{A + B} &= \overline{A} \cdot \overline{B} & \text{(De Morganovy zákony)}
\end{align*}
\end{vzorcebox}

\textbf{Funkčně úplné systémy}\\
NAND nebo NOR samy o sobě dokáží realizovat libovolnou logickou funkci.
To znamená, že kdybychom měli k dispozici pouze hradla NAND (nebo pouze NOR),
dokážeme sestavit jakýkoliv logický obvod.

\begin{center}
\begin{tikzpicture}[scale=0.8,
  gate/.style={rectangle, draw=accent, fill=lightblue, rounded corners=2pt,
               minimum width=1.2cm, minimum height=0.5cm, font=\small}]
  % NAND jako NOT
  \node[gate] (n1) at (0,0) {NAND};
  \draw[>-] (-0.8,-0.1) -- (n1.west);
  \draw[>-] (-0.8,0.1) -- (n1.west);
  \draw[>-] (n1.east) -- (0.8,0);
  \node[below, font=\scriptsize] at (0,-0.4) {NOT z NAND};

  % NAND jako AND
  \node[gate] (n2) at (2.5,0) {NAND};
  \node[gate] (n3) at (4.5,0) {NAND};
  \draw[>-] (1.7,-0.1) -- (n2.west);
  \draw[>-] (1.7,0.1) -- (n2.west);
  \draw[>-] (n2.east) -- (n3.west);
  \draw[>-] (n3.east) -- (n3.west);
  \draw[>-] (n3.east) -- (5.5,0);
  \node[below, font=\scriptsize] at (3.5,-0.4) {AND z NAND};
\end{tikzpicture}
\end{center}

\textbf{Programovací jazyky PLC dle normy IEC 61131-3}\\
Norma IEC 61131-3 definuje 5 programovacích jazyků pro PLC:

\begin{center}
\begin{tabularx}{0.95\linewidth}{@{} l l X @{}}
    \toprule
    \textbf{Zkratka} & \textbf{Název} & \textbf{Popis a použití} \\
    \midrule
    LD  & Ladder Diagram & Grafický, podobá se elektrickému schématu;
                           nejrozšířenější v průmyslu \\
    FBD & Function Block Diagram & Grafický, propojení funkčních bloků;
                                   vhodný pro regulaci \\
    ST  & Structured Text & Textový, podobný Pascalu;
                            vhodný pro složité algoritmy \\
    IL  & Instruction List & Nízkoúrovňový assembler PLC;
                             dnes již málo používaný \\
    SFC & Sequential Function Chart & Šablona stavového automatu;
                                       sekvenční procesy \\
    \bottomrule
\end{tabularx}
\end{center}

\textbf{Příklad kódu v jednotlivých jazycích:}

\textit{LD (Ladder Diagram):}
\begin{verbatim}
   I0.0      I0.1      Q0.0
|----| |------|/|------( )---|
|  start     stop     motor  |
\end{verbatim}

\textit{ST (Structured Text):}
\begin{verbatim}
IF start AND NOT stop THEN
    motor := TRUE;
ELSE
    motor := FALSE;
END_IF;
\end{verbatim}

\textit{SFC (Sequential Function Chart):}
\begin{verbatim}
[INIT] --(start)--> [RUN] --(stop)--> [STOP]
   ^                                   |
   +-----------------------------------+
\end{verbatim}

\textbf{Proč více jazyků?}\\
Různí výrobci PLC (Siemens, Allen-Bradley, Schneider) mají vlastní tradice.
Norma IEC 61131-3 umožňuje přenositelnost programů mezi platformami a usnadňuje
komunikaci mezi dodavateli.
\end{vysvetlenibox}

\newpage
\begin{vysvetlenibox}
	\textbf{Příklad}\\
	Jsou tři větráky a, b, c a alarm x, pokud dva a více větráků nefungují, spustí se alarm
	
	\textbf{Pravdivostní tabulka}
	\[
	\begin{array}{|c|c|c||c|}
		\hline
		A & B & C & X \\
		\hline
		0 & 0 & 0 & 1 \\
		0 & 0 & 1 & 1 \\
		0 & 1 & 0 & 1 \\
		0 & 1 & 1 & 0 \\
		\hline
		1 & 0 & 0 & 1 \\
		1 & 0 & 1 & 0 \\
		1 & 1 & 0 & 0 \\
		1 & 1 & 1 & 0 \\
		\hline

	\end{array}
	\]
	
\textbf{Karnaughova mapa}
\begin{center}
	\begin{karnaugh-map}[4][2][1][$BA$][$C$]
		\minterms{0, 1, 2, 4}
		\maxterms{3, 5, 6, 7}
		
		\implicant{0}{4}
		\implicant{0}{1}
		\implicantedge{0}{0}{2}{2}
		
	\end{karnaugh-map}
\end{center}

\textbf{Přepis}\\
	\[x = \neg a \cdot \neg b + \neg a \cdot \neg c + \neg b \cdot \neg c\]
	
\textbf{Schéma}\\

\begin{circuitikz}[scale=0.9, transform shape]
	
	% NOT hradla
	\node[not port] (nota) at (2, 6)  {};
	\node[not port] (notb) at (2, 3)  {};
	\node[not port] (notc) at (2, 0)  {};
	
	% Vstupy
	\draw (nota.in) -- ++(-1.5,0) node[left] {$a$};
	\draw (notb.in) -- ++(-1.5,0) node[left] {$b$};
	\draw (notc.in) -- ++(-1.5,0) node[left] {$c$};
	
	% AND hradla
	\node[and port] (and1) at (6.5, 5)  {};
	\node[and port] (and2) at (6.5, 3)  {};
	\node[and port] (and3) at (6.5, 1)  {};
	
	% Spoje NOT a -> AND1(in1), AND2(in1)
	\draw (nota.out) -- ++(0.4,0) coordinate (Ja);
	\draw (Ja) -- (Ja |- and1.in 1) -- (and1.in 1);
	\draw (Ja) -- (Ja |- and2.in 1) -- (and2.in 1);
	\filldraw (Ja) circle (1.5pt);
	
	% Spoje NOT b -> AND1(in2), AND3(in1)
	\draw (notb.out) -- ++(0.2,0) coordinate (Jb);
	\draw (Jb) -- (Jb |- and1.in 2) -- (and1.in 2);
	\draw (Jb) -- (Jb |- and3.in 1) -- (and3.in 1);
	\filldraw (Jb) circle (1.5pt);
	
	% Spoje NOT c -> AND2(in2), AND3(in2)
	\draw (notc.out) -- ++(0.0,0) coordinate (Jc);
	\draw (Jc) -- (Jc |- and2.in 2) -- (and2.in 2);
	\draw (Jc) -- (Jc |- and3.in 2) -- (and3.in 2);
	\filldraw (Jc) circle (1.5pt);
	
	% OR hradla
	\node[or port] (or1) at (10.5, 4) {};
	\node[or port] (or2) at (14,   2.5) {};
	
	% AND -> OR1
	\draw (and1.out) -- (and1.out |- or1.in 1) -- (or1.in 1);
	\draw (and2.out) -- (and2.out |- or1.in 2) -- (or1.in 2);
	
	% OR1 + AND3 -> OR2
	\draw (or1.out)  -- (or1.out  |- or2.in 1) -- (or2.in 1);
	\draw (and3.out) -- (and3.out |- or2.in 2) -- (or2.in 2);
	
	% Vystup
	\draw (or2.out) -- ++(1,0) node[right] {$x$};
	
	% Popisky
	\node[font=\small, above=3pt] at (and1.north) {$\lnot a \cdot \lnot b$};
	\node[font=\small, above=3pt] at (and2.north) {$\lnot a \cdot \lnot c$};
	\node[font=\small, below=3pt] at (and3.south) {$\lnot b \cdot \lnot c$};
	
\end{circuitikz}
\end{vysvetlenibox}

% ============================================================

% ============================================================

\newpage

% ============================================================
\section{Kombinační logické řízení, minimalizace}

\begin{pojmybox}
\textbf{Klíčové pojmy:}
kombinační obvod, DNF, KNF, Karnaughova mapa, minterm, maxterm,
Quine-McCluskey, multiplexer, demultiplexer, dekodér, enkodér, sčítačka
\end{pojmybox}

\begin{vysvetlenibox}
\textbf{Kombinační obvod}\\
Výstup závisí \textit{pouze} na okamžitých vstupech -- obvod nemá paměť.
Logická funkce se popisuje pravdivostní tabulkou nebo algebraickým výrazem.

\textbf{Příklad: Hlasovací systém}\\
Mějme 3 členy komise (A, B, C). Návrh projde, když hlasují \textit{alespoň 2} pro.
\begin{center}
\begin{tabular}{ccc|c}
    \toprule
    A & B & C & Výstup Y \\
    \midrule
    0 & 0 & 0 & 0 \\
    0 & 0 & 1 & 0 \\
    0 & 1 & 0 & 0 \\
    0 & 1 & 1 & 1 \\
    1 & 0 & 0 & 0 \\
    1 & 0 & 1 & 1 \\
    1 & 1 & 0 & 1 \\
    1 & 1 & 1 & 1 \\
    \bottomrule
\end{tabular}
\end{center}

\textbf{DNF a KNF}\\
Postup vytvoření normálních forem z pravdivostní tabulky:
\begin{itemize}[noitemsep, leftmargin=1.4em, topsep=4pt]
    \item \textbf{DNF} (Disjunktivní normální forma) -- součet součinů
          (mintermů); pro každý řádek, kde výstup = 1.
          \textit{Příklad:} $Y = \overline{A}BC + A\overline{B}C + AB\overline{C} + ABC$
    \item \textbf{KNF} (Konjunktivní normální forma) -- součin součtů
          (maxtermů); pro každý řádek, kde výstup = 0.
          \textit{Příklad:} $Y = (A+B+C)(A+B+\overline{C})(A+\overline{B}+C)(\overline{A}+B+C)$
\end{itemize}

\textbf{Postup minimalizace pomocí Karnaughovy mapy:}
\begin{enumerate}[noitemsep, leftmargin=1.4em, topsep=4pt]
    \item Sestrojíme K-mapu z pravdivostní tabulky
    \item Seskupíme jedničky do skupin o velikosti $2^n$ (1, 2, 4, 8\ldots)
    \item Skupiny mohou přesahovat přes okraje (mapa je "srolovaná")
    \item Z každé skupiny vyvodíme zjednodušený výraz
    \item Výsledná funkce je součtem všech skupin
\end{enumerate}

\textbf{Karnaughova mapa -- postup minimalizace}\\
Grafická metoda pro minimalizaci až 4 proměnných.
Sousední buňky se liší v \textit{jedné} proměnné (Grayův kód).
Seskupujeme jedničky do skupin o velikosti $2^n$ (1, 2, 4, 8\ldots).
Čím větší skupina, tím jednodušší výsledný výraz.

\textbf{Pravidla pro Karnaughovu mapu:}
\begin{itemize}[noitemsep, leftmargin=1.4em, topsep=4pt]
    \item \textbf{Počet buněk:} $2^n$ kde $n$ je počet proměnných (2→4, 3→8, 4→16)
    \item \textbf{Číslování os:} Grayův kód (00, 01, 11, 10) -- sousedi se liší o 1 bit
    \item \textbf{Seskupování:} skupiny musí být obdélníkové, velikost $2^n$
    \item \textbf{Přesahy:} mapa je "nekonečná" -- levý okraj sousedí s pravým
    \item \textbf{Cíl:} co nejméně skupin, každá co největší
\end{itemize}

\begin{center}
\begin{tikzpicture}[scale=0.75]
  % Karnaughova mapa 3 proměnné (A, BC)
  \draw[accent, thick] (0,0) grid (4,2);
  % Záhlaví
  \node at (-0.6, 1.5) {\small $A=0$};
  \node at (-0.6, 0.5) {\small $A=1$};
  \node at (0.5, 2.3)  {\small $BC{=}00$};
  \node at (1.5, 2.3)  {\small $BC{=}01$};
  \node at (2.5, 2.3)  {\small $BC{=}11$};
  \node at (3.5, 2.3)  {\small $BC{=}10$};
  % Hodnoty (příklad: f = A·B + C)
  \node at (0.5,1.5) {0};
  \node at (1.5,1.5) {1};
  \node at (2.5,1.5) {1};
  \node at (3.5,1.5) {0};
  \node at (0.5,0.5) {0};
  \node at (1.5,0.5) {1};
  \node at (2.5,0.5) {1};
  \node at (3.5,0.5) {1};
  % Skupiny
  \draw[stripeblue, very thick, rounded corners=4pt]
    (1.08,0.08) rectangle (2.92,1.92);
  \draw[red!70!black, very thick, rounded corners=4pt]
    (2.08,-0.05) rectangle (4.05,0.95);
  % Popisky skupin
  \node[color=stripeblue, font=\scriptsize] at (2,2.7) {skupina: $B$};
  \node[color=red!70!black, font=\scriptsize] at (3.5,-0.4) {skupina: $A\cdot C$};
\end{tikzpicture}
\end{center}

\begin{vzorcebox}
Výsledek minimalizace: $f = B + A \cdot \overline{C}$
\quad (místo původního rozsáhlého DNF výrazu)
\end{vzorcebox}

\textbf{Quine-McCluskeyho metoda}\\
Pro více než 4 proměnné se používá algoritmická metoda:
\begin{enumerate}[noitemsep, leftmargin=1.4em, topsep=4pt]
    \item Seřadíme mintermy podle počtu jedniček v binárním zápisu
    \item Hledáme dvojice, které se liší v jednom bitu a slučujeme je
    \item Opakujeme, dokud nelze dále slučovat
    \item Z výsledných prvočlenů vybereme minimální pokrytí
\end{enumerate}
Tato metoda je vhodná pro počítačovou implementaci.

\textbf{Typické kombinační obvody}

\begin{center}
\begin{tabularx}{0.95\linewidth}{@{} l X @{}}
    \toprule
    \textbf{Obvod}    & \textbf{Funkce} \\
    \midrule
    Multiplexer (MUX) & Vybírá 1 ze $2^n$ vstupů a přepošle na výstup dle adresních vodičů \\
    Demultiplexer     & Směruje 1 vstup na 1 z $2^n$ výstupů \\
    Dekodér           & Převádí $n$-bitový kód na $2^n$ výstupů (aktivní vždy jen 1) \\
    Enkodér           & Opak dekodéru -- $2^n$ vstupů na $n$-bitový kód \\
    Polosčítačka      & Sčítá 2 bity, výstupy: součet $S$ a přenos $C$ \\
    Úplná sčítačka    & Sčítá 2 bity + vstupní přenos; základ ALU \\
    Komparátor        & Porovnává dvě $n$-bitová čísla ($A>B$, $A=B$, $A<B$) \\
    \bottomrule
\end{tabularx}
\end{center}

\textbf{Multiplexer v praxi}\\
Multiplexer funguje jako "číselný přepínač". Příklad: 4:1 MUX má 4 datové vstupy
($D_0$--$D_3$), 2 adresní vstupy ($A_0$, $A_1$) a 1 výstup ($Y$).
\begin{itemize}[noitemsep, leftmargin=1.4em, topsep=4pt]
    \item Pokud $A_1A_0 = 00$, na výstup jde $D_0$
    \item Pokud $A_1A_0 = 01$, na výstup jde $D_1$
    \item Pokud $A_1A_0 = 10$, na výstup jde $D_2$
    \item Pokud $A_1A_0 = 11$, na výstup jde $D_3$
\end{itemize}
\textit{Použití:} výběr signálu z více zdrojů, paralelně-sériový převod.

\textbf{Dekodér a enkodér}\\
\begin{itemize}[noitemsep, leftmargin=1.4em, topsep=4pt]
    \item \textbf{Dekodér 3:8} -- 3 bity na vstupu aktivují 1 z 8 výstupů.
          Použití: adresa paměti, výběr periférie, řízení 7-segmentového displeje.
    \item \textbf{Enkodér 8:3} -- 8 vstupů, aktivní vždy jen jeden; výstup je 3bitový kód.
          Prioritní enkodér řeší situaci, kdy je aktivních více vstupů (vybere nejvyšší).
\end{itemize}

\textbf{Sčítačky a ALU}\\
Polosčítačka sčítá 2 bity, ale neumí přenos ze sčítání nižšího řádu.
\begin{itemize}[noitemsep, leftmargin=1.4em, topsep=4pt]
    \item \textbf{Polosčítačka} -- 2 vstupy ($A$, $B$), 2 výstupy ($S$, $C_{out}$)
    \item \textbf{Úplná sčítačka} -- 3 vstupy ($A$, $B$, $C_{in}$), 2 výstupy ($S$, $C_{out}$)
    \item \textbf{4bitová sčítačka} -- 4 úplné sčítačky za sebou; $C_{out}$ z předchozí jde do $C_{in}$ další
\end{itemize}
ALU (Arithmetic Logic Unit) v procesoru obsahuje sčítačky, logické funkce
a posuvné registry -- vše jsou kombinační obvody.

\begin{vzorcebox}
Polosčítačka: $S = A \oplus B$, \quad $C = A \cdot B$\\
Úplná sčítačka: $S = A \oplus B \oplus C_{in}$, \quad
$C_{out} = A \cdot B + C_{in} \cdot (A \oplus B)$\\
4bitová sčítačka: $C_{out}$ z bitu 0 jde do $C_{in}$ bitu 1 atd.
\end{vzorcebox}

\textbf{Komparátor}\\
Porovnává dvě čísla $A$ a $B$. Výstupy: $A>B$, $A=B$, $A<B$.
\begin{itemize}[noitemsep, leftmargin=1.4em, topsep=4pt]
    \item Rovnost: $A=B$ když všechny bity souhlasí ($XNOR$ každého páru bitů)
    \item Větší: $A>B$ když nejvyšší neshodný bit je u $A$ roven 1
\end{itemize}
\textit{Použití:} řízení procesů, bezpečnostní limity, třídění dat.
\end{vysvetlenibox}

% ============================================================

% ============================================================

\newpage

% ============================================================
\section{Sekvenční logické řízení, klopné obvody}

\begin{pojmybox}
\textbf{Klíčové pojmy:}
sekvenční obvod, klopný obvod, SR, JK, D, T,
automat (Mealy, Moore), stavový diagram, čítač, registr,
synchronní/asynchronní řízení, hazard
\end{pojmybox}

\begin{vysvetlenibox}
\textbf{Sekvenční obvod}\\
Na rozdíl od kombinačního má sekvenční obvod \textbf{paměť} -- výstup závisí
na vstupech \textit{a} na předchozím stavu. Základním stavebním prvkem je
\textbf{klopný obvod} (bistabilní člen -- může být ve stavu 0 nebo 1).

\textbf{Rozdělení sekvenčních obvodů:}
\begin{itemize}[noitemsep, leftmargin=1.4em, topsep=4pt]
    \item \textbf{Asynchronní} -- stav se mění ihned při změně vstupů;
          rychlejší, ale náchylné k hazardům
    \item \textbf{Synchronní} -- stav se mění pouze při hodinovém impulsu (CLK);
          stabilnější, snadněji se navrhují
\end{itemize}

\textbf{Princip klopného obvodu}\\
Klopný obvod má dva stabilní stavy (0 a 1). Přechod mezi stavy se řídí:
\begin{itemize}[noitemsep, leftmargin=1.4em, topsep=4pt]
    \item \textbf{Úrovní} -- obvod reaguje, když je CLK na vysoké/nízké úrovni
    \item \textbf{Hranou} -- obvod reaguje pouze při náběžné/sestupné hraně CLK
\end{itemize}
Hranové řízení je výhodnější -- zabraňuje vícečetným přepnutím.

\textbf{Přehled klopných obvodů}

\begin{center}
\begin{tabularx}{0.95\linewidth}{@{} l l X l @{}}
    \toprule
    \textbf{Typ} & \textbf{Vstupy} & \textbf{Funkce} & \textbf{Zvláštnost} \\
    \midrule
    RS  & S, R       & S=1: nastav Q=1; R=1: nastav Q=0
                     & Zakázaný stav S=R=1 \\
    D   & D, CLK     & Q přejde na hodnotu D při hraně CLK
                     & Jednoduchý, základ registrů \\
    JK  & J, K, CLK  & J=K=0: paměť; J=1,K=0: set; J=0,K=1: reset;
                       J=K=1: přepnutí
                     & Nemá zakázaný stav \\
    T   & T, CLK     & T=1: přepne stav; T=0: paměť
                     & Základ čítačů \\
    \bottomrule
\end{tabularx}
\end{center}

\textbf{Detailní popis klopných obvodů:}
\begin{itemize}[noitemsep, leftmargin=1.4em, topsep=4pt]
    \item \textbf{RS klopný obvod} -- (set-reset) nejzákladnější typ; S (Set) nastaví Q=1,
          R (Reset) nastaví Q=0. Stav S=R=1 $\rightarrow$ zakázaný stav.
    \item \textbf{D klopný obvod} -- eliminuje zakázaný stav; data (D) se
          zapíší při hraně hodinového signálu (CLK). Základ registrů a pamětí.
    \item \textbf{JK klopný obvod} -- vylepšená verze RS obvodu.
    \begin{center}
    \begin{tabular}{|c|c|l|}
    	\hline
    	J & K & Výtup\\
    	\hline
    	0 & 0 & drží stav\\
    	0 & 1 & vynuluje\\
    	1 & 0 & nastaví\\
    	1 & 1 & překlopí\\
    	\hline
    \end{tabular}
	\end{center}
    \item \textbf{T klopný obvod} -- zjednodušený JK s jedním vstupem;
          při T=1 a CLK=1, se stav překlopí.
\end{itemize}

\begin{vzorcebox}
JK klopný obvod -- charakteristická rovnice:
$Q_{n+1} = J \cdot \overline{Q_n} + \overline{K} \cdot Q_n$\\
D klopný obvod: $Q_{n+1} = D$\\
T klopný obvod: $Q_{n+1} = T \oplus Q_n$\\
SR klopný obvod: $Q_{n+1} = S + \overline{R} \cdot Q_n$ (pro S·R=0)
\end{vzorcebox}

\textbf{Stavový diagram -- vizualizace chování}\\
Stavový diagram znázorňuje přechody mezi stavy pomocí kruhů a šipek:
\begin{itemize}[noitemsep, leftmargin=1.4em, topsep=4pt]
    \item \textbf{Kruh} -- reprezentuje stav (0 nebo 1)
    \item \textbf{Šipka} -- přechod mezi stavy při splnění podmínky
    \item \textbf{Podmínka} -- zápis ve tvaru "vstup/výstup"
\end{itemize}

\textbf{Moorův a Mealyho automat}\\
Automat je model sekvenčního obvodu s definovanými stavy a přechody:

\begin{center}
\begin{tabularx}{0.95\linewidth}{@{} l X X @{}}
    \toprule
    & \textbf{Moorův automat} & \textbf{Mealyho automat} \\
    \midrule
    Výstup závisí na & pouze na stavu & stavu i vstupu \\
    Stabilita výstupu & stabilní, mění se při přechodu & může se měnit i bez přechodu \\
    Počet stavů & obvykle více & obvykle méně \\
    Použití & bezpečnostní systémy & rychlé řídicí automaty \\
    \bottomrule
\end{tabularx}
\end{center}

\textbf{Příklad: Automat pro rozpoznání sekvence 101}\\
\begin{itemize}[noitemsep, leftmargin=1.4em, topsep=4pt]
    \item \textbf{Stavy:} S0 (čekám na 1), S1 (čekám na 0), S2 (čekám na 1), S3 (detekováno)
    \item \textbf{Přechody:} S0--1→S1, S1--0→S2, S2--1→S3, S3--0→S2, S3--1→S1
    \item \textbf{Výstup:} Y=1 pouze ve stavu S3 (Mealyho varianta)
\end{itemize}

\textbf{Čítače a registry}

\begin{itemize}[noitemsep, leftmargin=1.4em, topsep=4pt]
    \item \textbf{Asynchronní čítač} -- výstup předchozího KO je CLK
          následujícího; jednodušší, ale pomalejší (ripple carry).
          \textit{Nevýhoda:} zpoždění se sčítá při průchodu řetězcem.
    \item \textbf{Synchronní čítač} -- všechny KO mají společný CLK;
          rychlejší, bez hazardů. Převodní logika určuje, které KO se přepne.
    \item \textbf{Čítač se zátěží} -- lze přednastavit počáteční hodnotu;
          používá se pro děličky s programovatelným poměrem.
    \item \textbf{Posuvný registr} -- D klopné obvody zapojené do série;
          sériový vstup / výstup dat. Typy:
          \begin{itemize}[noitemsep, leftmargin=1em]
              \item \textbf{SIPO} (Serial-In Parallel-Out) -- sériový vstup, paralelní výstup
              \item \textbf{PISO} (Parallel-In Serial-Out) -- paralelní vstup, sériový výstup
              \item \textbf{SISO} (Serial-In Serial-Out) -- zpoždění signálu
              \item \textbf{PIPO} (Parallel-In Parallel-Out) -- přesun dat mezi registry
          \end{itemize}
\end{itemize}

\textbf{Hazardy v sekvenčních obvodech}\\
Hazard je krátký nežádoucí impuls při změně vstupů:
\begin{itemize}[noitemsep, leftmargin=1.4em, topsep=4pt]
    \item \textbf{Statický hazard} -- výstup má krátce opačnou hodnotu
    \item \textbf{Dynamický hazard} -- výstup několikanásobně přepne
    \item \textbf{Řešení:} synchronní design, redundance v logice, filtry
\end{itemize}
\end{vysvetlenibox}
% ============================================================

% ============================================================

\newpage

% ============================================================
\section{Diskrétní řízení -- použití, signály, diskretizace, PSD regulátor}

\begin{pojmybox}
\textbf{Klíčové pojmy:}
diskrétní signál, vzorkování, Shannonův teorém, kvantování,
A/D převodník, perioda vzorkování $T_s$, Z-transformace,
PSD regulátor, dopředná/zpětná diference, antiwindup
\end{pojmybox}

\begin{vysvetlenibox}
\textbf{Proč diskrétní řízení?}\\
Moderní řídicí systémy (PLC, MCU) jsou digitální -- zpracovávají signál
\textit{nespojitě}, v pravidelných časových okamžicích.
Spojitý signál musíme nejprve \textbf{vzorkovat} a \textbf{kvantovat}.

\textbf{Proces převodu spojitého signálu na diskrétní:}
\begin{enumerate}[noitemsep, leftmargin=1.4em, topsep=4pt]
    \item \textbf{Vzorkování} -- získání hodnoty signálu v časových okamžicích $t = k \cdot T_s$
    \item \textbf{Držení} -- hodnota se udržuje konstantní po dobu převodu (sample \& hold)
    \item \textbf{Kvantování} -- zaokrouhlení napětí na nejbližší diskrétní úroveň
    \item \textbf{Kódování} -- přiřazení binárního kódu kvantizační úrovni
\end{enumerate}

\textbf{Shannonův vzorkovací teorém}

\begin{vzorcebox}
$$f_s \geq 2 \cdot f_{max}$$
Vzorkovací frekvence musí být alespoň \textbf{dvojnásobek} nejvyšší
frekvence v signálu. Jinak dojde k \textbf{aliasingu} (překrývání spekter).
\end{vzorcebox}

\textbf{Příklad aliasingu:}\\
Pokud vzorkujeme signál 50\,Hz frekvencí 80\,Hz ($< 2 \cdot 50$),
výsledek bude vypadat jako signál s frekvencí $80 - 50 = 30$\,Hz.
Tento falešný signál se nazývá \textit{alias}.

\textbf{Antialiasingový filtr}\\
Abychom předešli aliasingu, řadíme před A/D převodník \textbf{dolní propust},
k která potlačí frekvence vyšší než $f_s / 2$. Typicky se používá
Besselův nebo Butterworthův filtr 2.--4. řádu.

\textbf{Kvantování a rozlišení A/D převodníku}

\begin{vzorcebox}
$$\Delta U = \frac{U_{ref}}{2^n - 1}$$
kde $n$ je počet bitů A/D převodníku, $\Delta U$ je kvantizační krok.\\
Příklad: 12-bitový ADC s $U_{ref} = 3{,}3\,\text{V}$:
$\Delta U = \frac{3{,}3}{4095} \approx 0{,}806\,\text{mV}$
\end{vzorcebox}

\textbf{Typy A/D převodníků:}
\begin{itemize}[noitemsep, leftmargin=1.4em, topsep=4pt]
    \item \textbf{Flash ADC} -- nejrychlejší (ns), mnoho komparátorů,
          pro vysoké frekvence (osciloskopy, rádiové přijímače)
    \item \textbf{SAR ADC} (Successive Approximation) -- střední rychlost ($\mu$s),
          nízká spotřeba; nejpoužívanější v mikrokontrolérech
    \item \textbf{Sigma-Delta ADC} -- vysoké rozlišení (24 bitů), pomalejší;
          pro audio, přesná měření (vážní snímače, teploměry)
    \item \textbf{Dual-slope ADC} -- velmi přesný, pomalý;
          pro multimetry a měřicí přístroje
\end{itemize}

\textbf{Parametry A/D převodníků:}
\begin{itemize}[noitemsep, leftmargin=1.4em, topsep=4pt]
    \item \textbf{Rozlišení} -- počet bitů ($n$); určuje počet úrovní $2^n$
          (8 bitů = 256, 10 bitů = 1024, 12 bitů = 4096, 16 bitů = 65\,536)
    \item \textbf{Vzorkovací frekvence} -- maximální počet převodů za sekotu
          (kS/s, MS/s); např. Arduino Uno: 10-bitový ADC, 15\,kS/s
    \item \textbf{Referenční napětí $U_{ref}$} -- určuje rozsah měření;
          může být interní (např. 1,1\,V, 2,5\,V) nebo externí
    \item \textbf{Kvantizační chyba} -- $\pm \frac{1}{2}$ LSB (Least Significant Bit);
          nelze odstranit, je vlastností převodu
    \item \textbf{Nelinearita} -- INL (Integral Non-Linearity) a DNL (Differential);
          udává odchylku od ideální přímky
    \item \textbf{Čas převodu} -- doba mezi startem převodu a platným výstupem
\end{itemize}

\textbf{Zapojení A/D převodníku v obvodu:}
\begin{itemize}[noitemsep, leftmargin=1.4em, topsep=4pt]
    \item \textbf{Vstupní impedance} -- ADC zatěžuje zdroj; vysoká impedance zdroje
          způsobuje chyby (doporučuje se buffer/OP)
    \item \textbf{Filtr na vstupu} -- RC dolní propust před ADC potlačí vysoké frekvence
          a sníží aliasing i šum
    \item \textbf{Stínění} -- analogové vedení vést odděleně od digitálních signálů;
          zemní plocha jako stínění
\end{itemize}

\textbf{D/A převodníky}\\
Pro převod digitálního signálu na analogový (např. řízení motoru, generování signálů):

\textbf{Princip funkce:}
\begin{itemize}[noitemsep, leftmargin=1.4em, topsep=4pt]
    \item \textbf{R-2R žebřík} -- rezistorová síť, každý bit přispívá proudem
          podle své váhy; MSB má největší vliv, LSB nejmenší
    \item \textbf{PWM + filtr} -- pulzně šířková modulace s RC filtrem;
          jednoduché, levné, pro nenáročné aplikace (stmívání LED, řízení otáček)
    \item \textbf{DAC s proudovým výstupem} -- přesné zdroje proudu řízené bity;
          vyžaduje převodník proudu na napětí (OP)
\end{itemize}

\textbf{Parametry D/A převodníků:}
\begin{itemize}[noitemsep, leftmargin=1.4em, topsep=4pt]
    \item \textbf{Rozlišení} -- počet bitů; určuje minimální změnu výstupního napětí
    \item \textbf{Usazovací čas (settling time)} -- doba, za kterou výstup dosáhne
          cílové hodnoty s danou přesností (typicky 1--10\,$\mu$s)
    \item \textbf{Glitch} -- přechodový jev při změně kódu (např. 0111 $\to$ 1000);
          potlačuje se sample\&hold obvodem na výstupu
    \item \textbf{Reference} -- interní nebo externí; určuje rozsah výstupního napětí
\end{itemize}

\textbf{Použití D/A převodníků:}
\begin{itemize}[noitemsep, leftmargin=1.4em, topsep=4pt]
    \item \textbf{Generování signálů} -- sinus, trojúhelník, pila pro audio,
          testovací generátory
    \item \textbf{Řízení pohonů} -- analogový řídicí signál pro frekvenční měniče,
          servozesilovače
    \item \textbf{Zdroje napětí} -- programovatelné laboratorní zdroje,
          kalibrátory měřicích přístrojů
\end{itemize}

\textbf{PSD regulátor -- diskrétní analogie PID}\\
Derivace se nahradí \textbf{diferencí}, integrál \textbf{sumou}:

\begin{vzorcebox}
$$u_k = K_p \cdot e_k
+ K_i \cdot T_s \sum_{j=0}^{k} e_j
+ K_d \cdot \frac{e_k - e_{k-1}}{T_s}$$
kde $e_k = w_k - y_k$ je regulační odchylka v kroku $k$,
$T_s$ je perioda vzorkování.
\end{vzorcebox}

\textbf{Implementační detaily PSD regulátoru:}
\begin{itemize}[noitemsep, leftmargin=1.4em, topsep=4pt]
    \item \textbf{Proporcionální složka} -- okamžitá reakce na chybu
    \item \textbf{Sumační člen} -- akumulace chyby; pozor na přetečení
    \item \textbf{Derivační člen} -- predikce budoucího vývoje;
          v praxi se často filtruje (šumová citlivost)
    \item \textbf{Pozicionální algoritmus} -- výpočet absolutní hodnoty $u_k$
    \item \textbf{Inkrementální algoritmus} -- výpočet změny $\Delta u_k$;
          bezpečnější při přepnutí do manuálního režimu
\end{itemize}

\textbf{Vliv periody vzorkování $T_s$}

\begin{center}
\begin{tabularx}{0.95\linewidth}{@{} l X X @{}}
    \toprule
    & \textbf{Příliš velké $T_s$} & \textbf{Příliš malé $T_s$} \\
    \midrule
    Stabilita     & systém může kmitat nebo být nestabilní & dobrá \\
    Výpočetní zátěž & nízká & vysoká \\
    Aliasing      & hrozba & nehrozí \\
    Doporučení    & $T_s < \frac{T_{63\%}}{10}$
                  & zbytečné plýtvání výkonem \\
    \bottomrule
\end{tabularx}
\end{center}

\textbf{Volba periody vzorkování v praxi:}
\begin{itemize}[noitemsep, leftmargin=1.4em, topsep=4pt]
    \item \textbf{Teplotní soustavy} -- pomalé ($T_{63\%} \approx 10$\,s): $T_s = 1$\,s
    \item \textbf{Hladiny, tlaky} -- střední ($T_{63\%} \approx 1$\,s): $T_s = 100$\,ms
    \item \textbf{Pohony, poloha} -- rychlé ($T_{63\%} \approx 100$\,ms): $T_s = 10$\,ms
    \item \textbf{Elektronické zatížení} -- velmi rychlé: $T_s = 10$\,--\,100\,$\mu$s
\end{itemize}

\textbf{Antiwindup}\\
Antiwindup je ochrana proti přetečení integrační složky při saturaci
akčního zásahu. Když regulátor dosáhne maxima (např. 100\,\% výkonu),
integrace se zastaví -- jinak by se "naintegrovala" velká hodnota
a výstup by pozdě přestál přeběh.

\begin{itemize}[noitemsep, leftmargin=1.4em, topsep=4pt]
    \item \textbf{Podmíněná integrace} -- integrace běží, pouze když $|e_k| < \text{limit}$
    \item \textbf{Back-calculation} -- rozdíl mezi ideálním a saturovaným výstupem
          se vrací na vstup integrátoru se zpožděním
    \item \textbf{Clamping} -- zastavení integrace při dosažení limitu
\end{itemize}

\textbf{Z-transformace}\\
Pro analýzu diskrétních systémů se používá Z-transformace:
$$X(z) = \sum_{k=0}^{\infty} x_k \cdot z^{-k}$$
Převod mezi spojitým a diskrétním světem:
\begin{itemize}[noitemsep, leftmargin=1.4em, topsep=4pt]
    \item \textbf{Dopředná diference} -- $s \approx \frac{z-1}{T_s}$
    \item \textbf{Zpětná diference} -- $s \approx \frac{1-z^{-1}}{T_s}$
    \item \textbf{Tustinova metoda} -- $s \approx \frac{2}{T_s} \cdot \frac{z-1}{z+1}$;
          zachovává stabilitu, nejpoužívanější
\end{itemize}
\end{vysvetlenibox}

% ============================================================

% ============================================================

\newpage

% ============================================================
\section{Spojité řízení a regulace -- PID regulátor}

\begin{pojmybox}
\textbf{Klíčové pojmy:}
zpětná vazba, odchylka, P, I, D složka, proporcionální zesílení,
integrační čas, derivační čas, přeregulace, ustálená odchylka,
ladění (Ziegler-Nichols), přenosová funkce regulátoru
\end{pojmybox}

\begin{vysvetlenibox}
\textbf{Princip zpětnovazebního řízení}\\
Regulace s uzavřenou smyčkou porovnává skutečnou hodnotu $y$ s požadovanou
hodnotou $w$ a vypočítává akční zásah $u$ tak, aby se odchylka $e = w - y$
minimalizovala.

\begin{center}
\begin{tikzpicture}[scale=0.9,
  block/.style={rectangle, draw=accent, fill=lightblue,
                minimum width=1.8cm, minimum height=0.7cm,
                text centered, font=\small},
  sum/.style={circle, draw=accent, fill=lightblue,
              minimum size=0.5cm, font=\small}]
  \node[sum]   (sum)  at (1,0)   {$\Sigma$};
  \node[block] (reg)  at (3,0)   {PID};
  \node[block] (akt)  at (5.2,0) {Aktuátor};
  \node[block] (sus)  at (7.4,0) {Soustava};
  \node[block] (sen)  at (7.4,-1.5) {Senzor};
  \node (w) at (-0.5, 0) {$w$};
  \node (y) at (9.2,  0) {$y$};
  \draw[>-] (w)    -- (sum);
  \draw[>-] (sum)  -- (reg) node[midway,above,font=\scriptsize]{$e$};
  \draw[>-] (reg)  -- (akt) node[midway,above,font=\scriptsize]{$u$};
  \draw[>-] (akt)  -- (sus);
  \draw[>-] (sus)  -- (y);
  \draw[>-] (8.3,0) --  (sen);
  \draw[>-] (sen)   -- (sum) node[near end, left, font=\scriptsize]{$-$};
\end{tikzpicture}
\end{center}

\textbf{Typy regulací podle žádané hodnoty:}
\begin{itemize}[noitemsep, leftmargin=1.4em, topsep=4pt]
    \item \textbf{Stabilizace} -- $w = \text{konst.}$; udržování konstantní teploty, otáček
    \item \textbf{Sledování} -- $w = f(t)$; robot sleduje trajektorii, anténa sleduje družici
    \item \textbf{Programová regulace} -- $w$ se mění podle předem daného programu
          (např. pec: ohřev $\to$ výdrž $\to$ chlazení)
\end{itemize}

\textbf{Matematický popis PID regulátoru}

\begin{vzorcebox}
$$u(t) = K_p \left[ e(t) + \frac{1}{T_i}\int_0^t e(\tau)\,d\tau
+ T_d \frac{de(t)}{dt} \right]$$
Přenosová funkce: $G_R(s) = K_p \left(1 + \frac{1}{T_i s} + T_d s\right)$
\end{vzorcebox}

\textbf{Fyzikální význam parametrů:}
\begin{itemize}[noitemsep, leftmargin=1.4em, topsep=4pt]
    \item \textbf{$K_p$} (proporcionální zesílení) -- určuje sílu reakce na chybu
    \item \textbf{$T_i$} (integrační čas) -- doba, za kterou I-složka "dohoní" P-složku
    \item \textbf{$T_d$} (derivační čas) -- jak dopředu se díváme pro predikci
    \item \textbf{Alternativní zápis:} $K_i = K_p/T_i$, $K_d = K_p \cdot T_d$
\end{itemize}

\textbf{Vliv jednotlivých složek}

\begin{center}
\begin{tabularx}{0.95\linewidth}{@{} l X X X @{}}
    \toprule
    \textbf{Složka} & \textbf{Efekt} & \textbf{Zvyšování}
                    & \textbf{Nebezpečí} \\
    \midrule
    P & Rychlá reakce & Rychlejší odezva & Trvalá odchylka, kmity \\
    I & Odstraní trvale odchylku & Rychlejší eliminace & Přeregulace, nestabilita \\
    D & Tlumí kmitání & Méně překmitu & Zesiluje šum \\
    \bottomrule
\end{tabularx}
\end{center}

\textbf{Projevy jednotlivých složek v čase:}
\begin{itemize}[noitemsep, leftmargin=1.4em, topsep=4pt]
    \item \textbf{P-složka} -- reaguje okamžitě na přítomnou chybu;
          chyba $\to$ akce; problém: zůstává trvalá odchylka
    \item \textbf{I-složka} -- reaguje na minulost (akumuluje chybu);
          čím déle trvá chyba, tím větší akce; eliminuje trvalou odchylku
    \item \textbf{D-složka} -- reaguje na budoucnost (predikce);
          když chyba rychle roste, D-složka brzdí; tlumí překmit
\end{itemize}

\textbf{Problémy v praxi:}
\begin{itemize}[noitemsep, leftmargin=1.4em, topsep=4pt]
    \item \textbf{Derivační ráz} -- skok žádané hodnoty $\to$ D-složka "vyšílí";
          řeší se derivace podle výstupu $y$ místo chyby $e$
    \item \textbf{Windup} -- při saturaci se integrátor "nabije";
          řeší se antiwindup (podmíněná integrace, back-calculation)
    \item \textbf{Šum} -- D-složka zesiluje vysokofrekvenční šum;
          řeší se filtrací měřeného signálu
\end{itemize}

\textbf{Metoda Ziegler-Nichols (oscilační)}\\
Jedna z nejznámějších metod ladění PID regulátoru, vyvinutá v roce 1942.
Postup pro určení optimálních parametrů:

\begin{enumerate}[noitemsep, leftmargin=1.4em, topsep=4pt]
    \item Nastavit $T_i = \infty$, $T_d = 0$ (čistý P regulátor)
    \item Zvyšovat $K_p$ až do trvalých nehasících kmitů (hranice stability)
    \item Odečíst kritické zesílení $K_{kr}$ a kritickou periodu $T_{kr}$
    \item Vypočítat parametry dle tabulky:
\end{enumerate}

\begin{center}
\begin{tabular}{@{} l c c c @{}}
    \toprule
    \textbf{Typ} & $K_p$ & $T_i$ & $T_d$ \\
    \midrule
    P   & $0{,}50\,K_{kr}$ & $\infty$ & $0$ \\
    PI  & $0{,}45\,K_{kr}$ & $0{,}83\,T_{kr}$ & $0$ \\
    PID & $0{,}60\,K_{kr}$ & $0{,}50\,T_{kr}$ & $0{,}125\,T_{kr}$ \\
    \bottomrule
\end{tabular}
\end{center}

\textbf{Nevýhody Ziegler-Nichols:}
\begin{itemize}[noitemsep, leftmargin=1.4em, topsep=4pt]
    \item Systém se musí dostat do kmitů -- nevhodné pro některé provozy
    \item Výsledek je často příliš "agresivní" (velké překmity)
    \item Neplatí pro soustavy s velkým dopravním zpožděním
\end{itemize}

\textbf{Další metody ladění:}
\begin{itemize}[noitemsep, leftmargin=1.4em, topsep=4pt]
    \item \textbf{Chua-Broida} -- založená na přechodové charakteristice;
          změříme čas $T_{63\%}$ a $T_{28\%}$, vypočítáme $T_i$, $T_d$
    \item \textbf{Heuris-tické} -- $T_i \approx 4 \cdot T_{63\%}$, $T_d \approx T_{63\%} / 2$
    \item \textbf{Auto-tuning} -- moderní regulátory samy najdou parametry
          (Siemens, ABB, Honeywell mají vestavěné auto-tuning funkce)
    \item \textbf{Optimalizační} -- minimalizace kritéria (ISE, IAE, ITAE)
          pomocí numerických metod
\end{itemize}

\textbf{Kritéria kvality regulace:}
\begin{itemize}[noitemsep, leftmargin=1.4em, topsep=4pt]
    \item \textbf{ISE} (Integral of Squared Error) -- $\int e^2 dt$;
          penalizuje velké chyby
    \item \textbf{IAE} (Integral of Absolute Error) -- $\int |e| dt$;
          rovnoměrná penalizace
    \item \textbf{ITAE} (Integral of Time-weighted AE) -- $\int t \cdot |e| dt$;
          penalizuje chyby, které přetrvávají dlouho
\end{itemize}

\begin{center}
\begin{tikzpicture}
\begin{axis}[
  width=9cm, height=4.5cm,
  xlabel={Čas $t$ [s]},
  ylabel={Výstup $y(t)$},
  xmin=0, xmax=10, ymin=0, ymax=1.5,
  grid=major, grid style={dotted, medgray},
  legend pos=south east,
  legend style={font=\small},
  tick label style={font=\small},
  label style={font=\small},
  title style={font=\small\bfseries},
  title={Srovnání odezvy regulátorů na skok žádané hodnoty}
]
\addplot[accent, thick, domain=0:10, samples=80]
    {1 - exp(-1.5*x)*(cos(deg(2*x)) + 0.5*sin(deg(2*x)))};
\addlegendentry{PID}
\addplot[stripeblue, dashed, thick, domain=0:10, samples=80]
    {1 - exp(-0.8*x)};
\addlegendentry{PI}
\addplot[red!70!black, dotted, thick, domain=0:10, samples=80]
    {1 - exp(-0.5*x) - 0.1};
\addlegendentry{P (trvalá odch.)}
\addplot[medgray, domain=0:10] {1};
\end{axis}
\end{tikzpicture}
\end{center}

\textbf{Příklad z praxe: Regulace teploty v peci}\\
Máme pec s topným tělesem (0--230\,V) a teploměrem (Pt100).
\begin{itemize}[noitemsep, leftmargin=1.4em, topsep=4pt]
    \item \textbf{Krok 1:} Nastavíme čistý P režim ($T_i=\infty$, $T_d=0$)
    \item \textbf{Krok 2:} Zvyšujeme $K_p$ od 1 do 50; při $K_p=42$ začne systém
          kmitat s periodou $T_{kr} = 120$\,s
    \item \textbf{Krok 3:} Pro PID: $K_p = 0{,}6 \cdot 42 = 25{,}2$,
          $T_i = 0{,}5 \cdot 120 = 60$\,s, $T_d = 0{,}125 \cdot 120 = 15$\,s
    \item \textbf{Krok 4:} Ověříme funkčnost, případně doladíme pro menší překmit
\end{itemize}
\end{vysvetlenibox}

% ============================================================

% ============================================================

\newpage

% ============================================================
\section{Soustavy, statické, dynamické a frekvenční charakteristiky}

\begin{pojmybox}
\textbf{Klíčové pojmy:}
přenosová funkce, řádová soustava, přechodová charakteristika,
impulsní odezva, Bodeův diagram, frekvenční charakteristika,
fázová bezpečnost, amplitudová bezpečnost, póly, nuly
\end{pojmybox}

\begin{vysvetlenibox}
\textbf{Přenosová funkce}\\
Přenosová funkce je matematický model, který popisuje vstupně-výstupní
vztah lineární soustavy v komplexní oblasti. Používá se Laplaceova
transformace, která převádí diferenciální rovnice na algebraické rovnice.

\begin{vzorcebox}
$$G(s) = \frac{Y(s)}{U(s)} = \frac{b_m s^m + \cdots + b_0}
{a_n s^n + \cdots + a_0}$$
kde $s = j\omega$ je komplexní proměnná (Laplaceova transformace).\\
Soustava 1. řádu: $G(s) = \frac{K}{Ts + 1}$ \quad
(zesílení $K$, časová konstanta $T$)\\
Soustava 2. řádu: $G(s) = \frac{\omega_n^2}{s^2 + 2\zeta\omega_n s + \omega_n^2}$
\quad ($\zeta$ = poměrný útlum)
\end{vzorcebox}

\textbf{Póly a nuly přenosové funkce:}
\begin{itemize}[noitemsep, leftmargin=1.4em, topsep=4pt]
    \item \textbf{Nuly} -- kořeny čitatele; frekvence, při kterých je výstup nulový
    \item \textbf{Póly} -- kořeny jmenovatele; určují stabilitu a dynamiku
    \item \textbf{Stabilita:} všechny póly musí mít zápornou reálnou část
          (být v levé polorovině komplexní roviny)
\end{itemize}

\textbf{Typické soustavy podle řádu:}
\begin{itemize}[noitemsep, leftmargin=1.4em, topsep=4pt]
    \item \textbf{0. řád} -- $G(s) = K$; okamžitá odezva bez zpoždění
          (ideální zesilovač, dělič napětí)
    \item \textbf{1. řád} -- jeden pól; exponenciální náběh s časovou konstantou $T$
          (RC článek, ohřev s tělesem, nádrž s odtokem)
    \item \textbf{2. řád} -- dva póly; může kmitat při nízkém útlumu
          (pružina-hmota, RLC obvod, kyvadlo)
    \item \textbf{3. a vyšší řád} -- složené soustavy; často se aproximují
          na soustavu 1. nebo 2. řádu s dopravním zpožděním
\end{itemize}

\textbf{Dopravní zpoždění}\\
Některé systémy mají čisté zpoždění (materiál putuje pásem, signál se
zpracovává): $G(s) = G_0(s) \cdot e^{-Ls}$, kde $L$ je doba zpoždění.
Ve frekvenční oblasti: $e^{-j\omega L}$ -- fázový posun lineární s $\omega$.

\textbf{Charakteristiky soustav}

\begin{center}
\begin{tabularx}{0.95\linewidth}{@{} l X X @{}}
    \toprule
    \textbf{Charakteristika} & \textbf{Co popisuje} & \textbf{Jak se měří} \\
    \midrule
    Statická & Závislost výstupu na vstupu v ustáleném stavu;
               sklon = statické zesílení $K$
             & Nastavíme různé vstupy, čekáme na ustálení \\
    Přechodová & Odezva na jednotkový skok;
                 odečítáme $T$, $K$, přeregulaci $\sigma$
              & Skok vstupu, záznam výstupu \\
    Frekvenční & Bodeův diagram: $|G(j\omega)|$ a $\angle G(j\omega)$ vs. $\omega$
              & Harmonické buzení různých frekvencí \\
    Impulsní & Odezva na Diracův impuls; derivace přechodové
              & Krátký intenzivní zásah, záznam odezvy \\
    \bottomrule
\end{tabularx}
\end{center}

\textbf{Parametry přechodové charakteristiky:}
\begin{itemize}[noitemsep, leftmargin=1.4em, topsep=4pt]
    \item \textbf{$T_{63\%}$} -- čas, za který výstup dosáhne 63\% ustálené hodnoty
          (pro soustavu 1. řádu: $T_{63\%} = T$)
    \item \textbf{$T_{95\%}$} -- čas dosažení 95\% (pro 1. řád: $\approx 3T$)
    \item \textbf{Přeregulace $\sigma$} -- maximální překmit nad ustálenou hodnotu [\%]
    \item \textbf{Doba ustálení $T_u$} -- čas, kdy se výstup udrží v pásmu $\pm 5\,\%$
\end{itemize}

\textbf{Bodeův diagram -- princip}\\
Bodeův diagram se skládá ze dvou grafů:
\begin{itemize}[noitemsep, leftmargin=1.4em, topsep=4pt]
    \item \textbf{Amplitudová charakteristika} -- $|G(j\omega)|$ [dB] vs. $\omega$ [log]
    \item \textbf{Fázová charakteristika} -- $\angle G(j\omega)$ [°] vs. $\omega$ [log]
\end{itemize}

\textbf{Asymptotický průběh pro soustavu 1. řádu:}
\begin{itemize}[noitemsep, leftmargin=1.4em, topsep=4pt]
    \item Pro $\omega < \omega_c$ (lomová frekvence): vodorovná čára na 0\,dB
    \item Pro $\omega > \omega_c$: klesající čára se sklonem $-20$\,dB/dekádu
    \item Fáze: přechod od 0° k $-90°$ v okolí $\omega_c$
\end{itemize}

\begin{center}
\begin{tikzpicture}
\begin{semilogxaxis}[
  width=9cm, height=4cm,
  xlabel={Frekvence $\omega$ [rad/s]},
  ylabel={Zesílení [dB]},
  xmin=0.01, xmax=100,
  ymin=-40, ymax=5,
  grid=major, grid style={dotted,medgray},
  tick label style={font=\small},
  label style={font=\small},
  title={Bodeův diagram $G(s)=\frac{1}{s+1}$},
  title style={font=\small\bfseries}
]
\addplot[accent, thick, domain=0.01:100, samples=200]
    {-10*log10(1 + x^2)};
\addplot[stripeblue, dashed, domain=0.01:1] {0};
\addplot[stripeblue, dashed, domain=1:100]  {-20*(log10(x))};
\draw[red!70, dashed] (axis cs:1,-40) -- (axis cs:1,5)
    node[above, font=\scriptsize, color=red!70] {$\omega_c=1/T$};
\end{semilogxaxis}
\end{tikzpicture}
\end{center}

\textbf{Pravidla pro kreslení Bodeova diagramu:}
\begin{itemize}[noitemsep, leftmargin=1.4em, topsep=4pt]
    \item \textbf{Zesílení $K$} -- posune amplitudu o $20\log_{10}K$ dB
    \item \textbf{Pól v počátku} ($1/s$) -- sklon $-20$\,dB/dek, fáze $-90°$
    \item \textbf{Pól reálný} ($1/(Ts+1)$) -- lom při $\omega=1/T$, sklon $-20$\,dB/dek
    \item \textbf{Dvojice pólů} (2. řád) -- sklon $-40$\,dB/dek, fáze $-180°$
    \item \textbf{Nula} -- opačný efekt než pól (kladný sklon)
\end{itemize}

\textbf{Stabilitní kritéria z Bodeova diagramu}

\begin{vzorcebox}
\textbf{Fázová bezpečnost:} $\varphi_m = 180° + \angle G(j\omega_0)$
kde $\omega_0$ je frekvence, při níž $|G|=0\,\text{dB}$.
Požadujeme $\varphi_m > 30°$ (typicky $45°$--$60°$).\\[4pt]
\textbf{Amplitudová bezpečnost:} $A_m = -|G(j\omega_{\pi})|$
kde $\omega_{\pi}$ je frekvence, při níž $\angle G = -180°$.
Požadujeme $A_m > 6\,\text{dB}$.
\end{vzorcebox}

\textbf{Význam bezpečností:}
\begin{itemize}[noitemsep, leftmargin=1.4em, topsep=4pt]
    \item \textbf{Fázová bezpečnost} -- jak daleko jsme od vzniku kmitů;
          malá $\varphi_m$ $\to$ velké překmity, oscilace
    \item \textbf{Amplitudová bezpečnost} -- jak můžeme zvýšit zesílení,
          než systém ztratí stabilitu
    \item \textbf{Typické hodnoty:} $\varphi_m = 45°\text{--}60°$, $A_m = 6\text{--}12$\,dB
\end{itemize}

\textbf{Nyquistův diagram}\\
Alternativa k Bodeovu diagramu -- komplexní přenosová funkce
$G(j\omega)$ se vykreslí do jedné křivky v komplexní rovině.
\begin{itemize}[noitemsep, leftmargin=1.4em, topsep=4pt]
    \item Vzdálenost od bodu $[-1, j0]$ určuje stabilitu
    \item Kritério: křivka nesmí obklopit bod $[-1, j0]$
    \item Výhoda: vidíme amplitudu i fázi najednou
\end{itemize}

\textbf{Identifikace soustavy z měření}\\
Postup, když neznáme matematický model:
\begin{enumerate}[noitemsep, leftmargin=1.4em, topsep=4pt]
    \item Změříme přechodovou charakteristiku (skok vstupu)
    \item Odečteme $T_{63\%}$, $T_{28\%}$, přeregulaci $\sigma$
    \item Typ soustavy: bez přeregulace $\to$ 1. řád; s přeregulací $\to$ 2. řád
    \item Vypočítáme parametry: $T = T_{63\%}$, $K = y(\infty)/u_{skok}$
    \item Ověříme na frekvenční charakteristice
\end{enumerate}
\end{vysvetlenibox}

% ============================================================

% ============================================================


\newpage


% ============================================================
\section{Základy PLC a mikrokontrolérů}

\begin{pojmybox}
PLC, scan cyklus, I/O, CPU, ladder diagram,
mikrokontrolér, GPIO, UART, SPI, I2C,
interrupt, timer, PWM, register, firmware
\end{pojmybox}

% ============================================================
\begin{vysvetlenibox}

\textbf{PLC (Programmable Logic Controller)}\\
PLC je průmyslový řídicí systém navržený pro spolehlivý a dlouhodobý provoz.

\textbf{Hlavní vlastnosti:}
\begin{itemize}
\item deterministické řízení (stejný čas cyklu)
\item odolnost vůči rušení (EMC)
\item modulární struktura
\item jednoduchá údržba
\end{itemize}

% -----------------------------
\textbf{Scan cyklus PLC}

\begin{center}
\begin{tikzpicture}[
node/.style={rectangle, draw=accent, fill=lightblue,
rounded corners=3pt, minimum width=3.2cm, minimum height=0.6cm, font=\small}
]

\node[node] (a) {1. Čtení vstupů};
\node[node, below=of a] (b) {2. Vykonání programu};
\node[node, below=of b] (c) {3. Zápis výstupů};
\node[node, below=of c] (d) {4. Diagnostika};

\draw[->] (a)--(b);
\draw[->] (b)--(c);
\draw[->] (c)--(d);
\draw[->] (d.west) -- ++(-1,0) -- (a.west);

\node[right=0.3cm of b, font=\scriptsize] {typicky 1–10 ms};

\end{tikzpicture}
\end{center}

 PLC nepracuje „real-time per event“, ale cyklicky.

\textbf{Programovací jazyky PLC (IEC 61131-3)}

PLC se programují podle normy IEC 61131-3, která definuje několik jazyků:

\begin{itemize}
\item \textbf{LD (Ladder Diagram)} – žebříčkový diagram
\item \textbf{FBD (Function Block Diagram)} – blokové schéma
\item \textbf{ST (Structured Text)} – textový jazyk podobný Pascalu
\item \textbf{IL (Instruction List)} – instrukční seznam (historický)
\item \textbf{SFC (Sequential Function Chart)} – sekvenční grafy
\end{itemize}

% -----------------------------
\textbf{LD – Ladder Diagram}

Nejpoužívanější v průmyslu. Vypadá jako elektrické schéma relé.

\begin{itemize}
\item levá strana = vstupy
\item pravá strana = výstupy
\item logika AND/OR pomocí kontaktů
\end{itemize}

Výhoda: snadné pochopení pro elektrikáře.

% -----------------------------
\textbf{ST – Structured Text}

Textový jazyk podobný C/Pascalu:

\begin{vzorcebox}
\begin{verbatim}
IF temperature > 50 THEN
    motor := TRUE;
ELSE
    motor := FALSE;
END_IF;
\end{verbatim}
\end{vzorcebox}

Výhoda:
\begin{itemize}
\item složitější algoritmy
\item matematika
\item PID regulace
\end{itemize}

% -----------------------------
\textbf{FBD – Function Block Diagram}

Grafické bloky:

\begin{itemize}
\item AND / OR / NOT bloky
\item funkční bloky (timer, counter)
\item propojení signálů mezi bloky
\end{itemize}

Výhoda: přehlednost u řízení procesů.

% -----------------------------
\textbf{SFC – Sequential Function Chart}

Řídí sekvence kroků:

\begin{itemize}
\item kroky (steps)
\item přechody (transitions)
\item paralelní větvení
\end{itemize}

Použití:
\begin{itemize}
\item výrobní linky
\item automatické cykly
\end{itemize}

% -----------------------------
\textbf{IL – Instruction List (zastaralý)}

Nízká úroveň, podobná assembleru.

Dnes už se prakticky nepoužívá (deprecated).


% ============================================================
\textbf{Mikrokontrolér (MCU)}\\
Malý počítač na jednom čipu.

\textbf{Obsahuje:}
\begin{itemize}
\item CPU
\item RAM / Flash
\item GPIO piny
\item periferie (UART, SPI, I2C, ADC, timers)
\end{itemize}

% ============================================================
\textbf{PLC vs mikrokontrolér}

\begin{tabularx}{\linewidth}{l X X}
\toprule
 & PLC & Mikrokontrolér \\
\midrule
Cíl & průmysl & vestavěné systémy \\
Cena & vysoká & nízká \\
Odolnost & velmi vysoká & střední \\
Programování & IEC 61131-3 & C / C++ \\
Flexibilita & omezená & vysoká \\
Reakce & cyklická & event-driven \\
\bottomrule
\end{tabularx}

% ============================================================
\textbf{GPIO (mikrokontrolér)}

\begin{itemize}
\item vstup = čtení senzoru
\item výstup = řízení LED, motoru
\item alternativní funkce = UART/SPI/I2C/PWM
\end{itemize}

% ============================================================
\textbf{Interrupts (přerušení)}

Přerušení umožňuje okamžitou reakci na událost.

\begin{itemize}
\item tlačítko
\item časovač
\item UART příjem dat
\item změna signálu
\end{itemize}

 Výhoda: nemusí se čekat v loop()

% ============================================================
\textbf{Timers (časovače)}

Použití:
\begin{itemize}
\item měření času
\item PWM generace
\item periodické úlohy
\end{itemize}

\[
f_{out} = \frac{f_{cpu}}{prescaler}
\]

% ============================================================
\textbf{PWM}

\begin{vzorcebox}
\[
D = \frac{t_{on}}{T} \cdot 100\,\%
\]
\end{vzorcebox}

Použití:
\begin{itemize}
\item řízení motorů
\item LED stmívání
\item regulace výkonu
\end{itemize}

% ============================================================
\textbf{Komunikace MCU}

\begin{itemize}
\item UART – jednoduchá sériová linka
\item SPI – rychlá synchronní komunikace
\item I2C – více zařízení na 2 vodičích
\end{itemize}

% ============================================================
\textbf{Firmware}

Firmware = program v mikrokontroléru:
\begin{itemize}
\item běží přímo na HW
\item často bez OS
\item musí být stabilní
\end{itemize}

% ============================================================
\end{vysvetlenibox}



\newpage


% ============================================================
\section{Programování v jazyce C – základy}

\begin{pojmybox}
datové typy, proměnné, scope, local/global/static,
podmínky, cykly, funkce, pole
\end{pojmybox}

% ============================================================
\begin{vysvetlenibox}

\textbf{Proměnné a datové typy}

Proměnná je místo v paměti, které má jméno a typ.

\begin{vzorcebox}
\begin{verbatim}
int a = 10;
float b = 3.14;
char c = 'A';
\end{verbatim}
\end{vzorcebox}

\textbf{Rozdíl: local / global / static}

\begin{itemize}
\item \textbf{Local} – existuje jen uvnitř funkce
\item \textbf{Global} – existuje po celý program
\item \textbf{Static} – zachovává hodnotu mezi voláními funkce
\end{itemize}

\begin{vzorcebox}
\begin{verbatim}
int global = 5;   // globální

void f() {
    int local = 10;        // lokální
    static int count = 0;  // uchovává stav
    count++;
}
\end{verbatim}
\end{vzorcebox}

% ============================================================
\textbf{Podmínky (if / else)}

Slouží k rozhodování:

\begin{vzorcebox}
\begin{verbatim}
if (x > 10) {
    printf("větší než 10");
} else {
    printf("menší nebo rovno 10");
}
\end{verbatim}
\end{vzorcebox}

 Použití:
\begin{itemize}
\item řízení logiky programu
\item reakce na senzory
\item rozhodování v algoritmech
\end{itemize}

% ============================================================
\textbf{Cykly (loops)}

\textbf{For cyklus} – když známe počet opakování:

\begin{vzorcebox}
\begin{verbatim}
for (int i = 0; i < 10; i++) {
    printf("%d", i);
}
\end{verbatim}
\end{vzorcebox}

\textbf{While cyklus} – dokud platí podmínka:

\begin{vzorcebox}
\begin{verbatim}
while (x > 0) {
    x--;
}
\end{verbatim}
\end{vzorcebox}

\textbf{Do-while} – provede se alespoň jednou:

\begin{vzorcebox}
\begin{verbatim}
do {
    x--;
} while (x > 0);
\end{verbatim}
\end{vzorcebox}

% ============================================================
\textbf{Funkce}

Funkce slouží k opakovanému použití kódu:

\begin{vzorcebox}
\begin{verbatim}
int soucet(int a, int b) {
    return a + b;
}
\end{verbatim}
\end{vzorcebox}

 Výhody:
\begin{itemize}
\item přehlednost
\item opakované použití
\item modularita
\end{itemize}

% ============================================================
\textbf{Scope (viditelnost proměnných)}

\begin{itemize}
\item proměnná existuje jen v bloku { }
\item mimo blok ji nelze použít
\end{itemize}

\begin{vzorcebox}
\begin{verbatim}
{
    int x = 5;
}
// x zde už neexistuje
\end{verbatim}
\end{vzorcebox}

% ============================================================
\textbf{Shrnutí rozdílů}

\begin{tabularx}{\linewidth}{l X}
\toprule
\textbf{Typ} & \textbf{Popis} \\
\midrule
local & existuje uvnitř funkce \\
global & existuje celý program \\
static & uchovává hodnotu mezi voláními \\
\bottomrule
\end{tabularx}

% ============================================================
\end{vysvetlenibox}



\newpage


% ============================================================
\section{Průmyslové komunikace}

\begin{pojmybox}
sběrnice, RS-232, RS-485, Modbus, CAN,
SPI, I2C, I2S, UART,
master/slave, diferenciální přenos,
SCADA, OPC-UA
\end{pojmybox}

% ============================================================
\begin{vysvetlenibox}

\textbf{RS-232 vs RS-485}

\begin{center}
\begin{tabularx}{\linewidth}{l X X}
\toprule
 & RS-232 & RS-485 \\
\midrule
Typ přenosu & nesymetrický & diferenciální \\
Počet zařízení & 2 (bod-bod) & až 32+ zařízení \\
Dosah & cca 15 m & až 1200 m \\
Odolnost rušení & nízká & vysoká \\
Použití & PC komunikace & průmysl (Modbus RTU) \\
\bottomrule
\end{tabularx}
\end{center}

\textbf{Princip diferenciálního přenosu}

RS-485 používá dva vodiče (A, B):
\begin{itemize}
\item logická 1 = rozdíl napětí mezi vodiči
\item logická 0 = opačný rozdíl
\end{itemize}

Výhoda:
\begin{itemize}
\item vysoká odolnost proti rušení
\item vhodné pro dlouhé kabely
\end{itemize}

% ============================================================
\textbf{Průmyslové sběrnice}

\begin{tabularx}{\linewidth}{l X}
\toprule
Sběrnice & Použití \\
\midrule
Modbus & jednoduchá komunikace PLC \\
CAN & automobily, robustní komunikace \\
Profibus & starší Siemens systémy \\
Profinet & průmyslový Ethernet \\
EtherCAT & velmi rychlé řízení pohonů \\
\bottomrule
\end{tabularx}

% ============================================================
\textbf{Komunikace mikrokontrolérů}

\textbf{UART}
\begin{itemize}
\item asynchronní komunikace
\item 2 vodiče (TX, RX)
\item nutná shodná rychlost (baud rate)
\end{itemize}

\textbf{SPI}
\begin{itemize}
\item synchronní komunikace
\item vodiče: MOSI, MISO, SCK, CS
\item velmi vysoká rychlost
\item master řídí komunikaci
\end{itemize}

\textbf{I2C}
\begin{itemize}
\item 2 vodiče (SDA, SCL)
\item adresování zařízení
\item více zařízení na sběrnici
\item pomalejší než SPI
\end{itemize}

\textbf{I2S}
\begin{itemize}
\item přenos digitálního audia
\item používá se pro DAC/ADC
\item synchronní datový tok
\end{itemize}

% ============================================================
\textbf{Master/Slave komunikace}

\begin{itemize}
\item master řídí komunikaci
\item slave odpovídá na požadavky
\item příklad: Modbus
\end{itemize}

% ============================================================
\textbf{SCADA systém}

SCADA slouží k:
\begin{itemize}
\item vizualizaci dat
\item ovládání procesu
\item archivaci a diagnostice
\end{itemize}

% ============================================================
\textbf{OPC-UA}

Moderní komunikační standard:
\begin{itemize}
\item platformově nezávislý
\item zabezpečený (šifrování)
\item propojuje PLC, SCADA a IT systémy
\end{itemize}

% ============================================================
\end{vysvetlenibox}



\newpage


% ============================================================
\section{Senzorika -- poloha, rychlost, zrychlení, průtok, hladina}

\begin{pojmybox}
enkodér, LVDT, Hallův jev, MEMS, akcelerometr,
tachogenerátor, průtokoměr, hladinoměr,
citlivost, přesnost, rozlišení, linearita
\end{pojmybox}

% ============================================================
\begin{vysvetlenibox}

\textbf{Parametry senzorů}

\begin{vzorcebox}
\[
S = \frac{\Delta y}{\Delta x}
\]
\end{vzorcebox}

\begin{itemize}
\item \textbf{Citlivost} – změna výstupu na změnu vstupu
\item \textbf{Rozlišení} – nejmenší měřitelná změna
\item \textbf{Přesnost} – odchylka od skutečné hodnoty
\item \textbf{Linearita} – odchylka od ideální přímky
\end{itemize}

% ============================================================
\textbf{Snímače polohy}

\begin{tabularx}{\linewidth}{l X}
\toprule
Typ & Princip \\
\midrule
Inkrementální enkodér & počítá impulzy (relativní poloha) \\
Absolutní enkodér & každá poloha má unikátní kód \\
Potenciometr & změna odporu dle polohy \\
LVDT & změna indukčnosti (lineární posuv) \\
Hallův snímač & detekce magnetického pole \\
\bottomrule
\end{tabularx}

% ============================================================
\textbf{Rychlost a zrychlení}

\begin{vzorcebox}
\[
v = \frac{ds}{dt}, \quad a = \frac{dv}{dt}
\]
\end{vzorcebox}

\begin{itemize}
\item \textbf{Tachogenerátor} – napětí úměrné otáčkám:
\[
U = k \cdot n
\]

\item \textbf{Enkodér} – rychlost z impulzů:
\[
\omega = \frac{\Delta \varphi}{\Delta t}
\]

\item \textbf{Akcelerometr (MEMS)}
princip: malá hmotnost se při zrychlení posune → změna kapacity nebo odporu
\end{itemize}

% ============================================================
\textbf{Průtokoměry}

\begin{tabularx}{\linewidth}{l X}
\toprule
Typ & Princip \\
\midrule
Elektromagnetický & indukce napětí ve vodivé kapalině \\
Ultrazvukový & měření doby průchodu signálu \\
Vortex & vznik vírů za překážkou \\
Coriolis & měření hmotnostního průtoku (setrvačné síly) \\
\bottomrule
\end{tabularx}

% ============================================================
\textbf{Snímače hladiny}

\begin{tabularx}{\linewidth}{l X}
\toprule
Typ & Princip \\
\midrule
Plovákový & mechanický plovák kopíruje hladinu \\
Kapacitní & změna kapacity dle hladiny \\
Ultrazvukový & měření vzdálenosti od hladiny \\
Tlakový & hydrostatický tlak: $p = \rho g h$ \\
\bottomrule
\end{tabularx}

% ============================================================
\end{vysvetlenibox}



\newpage


\section{Senzorika -- teplota, síla, hmotnost, deformace, tlak}

\begin{pojmybox}
\textbf{Klíčové pojmy:}
termistor NTC/PTC, termočlánek, RTD (Pt100),
IR teploměr, tenzometr, Wheatstoneův můstek,
piezorezistivní senzor, piezoelektrický jev,
absolutní tlak, přetlak, podtlak
\end{pojmybox}

\begin{vysvetlenibox}

% ------------------------------------------------------------
\textbf{Teplotní senzory}

\begin{center}
\begin{tabularx}{0.95\linewidth}{l X l l}
Typ & Princip & Rozsah & Přesnost \\
\hline
Pt100 & Odpor platiny roste s teplotou & $-200$ až $+850$ °C & $\pm0{,}1$ °C \\
Termočlánek & Seebeckův jev (2 kovy) & $-270$ až $+2300$ °C & $\pm1$ °C \\
NTC & Odpor klesá s teplotou & $-50$ až $+150$ °C & $\pm0{,}5$ °C \\
IR senzor & Infračervené záření & $-50$ až $+2000$ °C & $\pm2$ °C \\
\end{tabularx}
\end{center}

\begin{vzorcebox}
Pt100: $R(T) = R_0 (1 + AT + BT^2)$

Seebeckův jev: $U = S \cdot (T_1 - T_2)$
\end{vzorcebox}

% ------------------------------------------------------------
\textbf{Síla, hmotnost a deformace}

\textbf{Základ:}
\begin{itemize}
\item Síla: $F = m \cdot g$
\item Hmotnost se měří nepřímo přes sílu (váhy)
\end{itemize}

\textbf{Tenzometr:}
\begin{itemize}
\item mění odpor při deformaci
\item nalepený na pružném materiálu
\end{itemize}

\begin{vzorcebox}
$GF = \frac{\Delta R / R}{\varepsilon}$

$\varepsilon = \frac{\Delta l}{l}$
\end{vzorcebox}

% ------------------------------------------------------------
\textbf{Wheatstoneův můstek}

Používá se pro přesné měření malých změn odporu:

\begin{vzorcebox}
$U_{out} = U_{in} \cdot \frac{R_1 R_3 - R_2 R_4}{(R_1+R_2)(R_3+R_4)}$
\end{vzorcebox}

\begin{itemize}
\item čtvrtmůstek -- 1 tenzometr
\item půlmůstek -- 2 tenzometry
\item plný můstek -- 4 tenzometry (nejvyšší citlivost)
\end{itemize}

% ------------------------------------------------------------
\textbf{Tlakové senzory}

\begin{center}
\begin{tabularx}{0.95\linewidth}{l X}
Typ & Princip \\
\hline
Piezorezistivní & deformace membrány → změna odporu \\
Kapacitní & změna vzdálenosti elektrod → změna kapacity \\
Piezoelektrický & generuje napětí při deformaci \\
\end{tabularx}
\end{center}

\textbf{Typy tlaku:}
\begin{itemize}
\item absolutní (vůči vakuu)
\item relativní (přetlak vůči atmosféře)
\item podtlak (nižší než atmosférický)
\end{itemize}

% ------------------------------------------------------------
\textbf{Shrnutí principů}

\begin{itemize}
\item Teplota → změna odporu nebo napětí
\item Síla → deformace → změna odporu
\item Tlak → deformace membrány
\item Hmotnost → síla (gravitační)
\end{itemize}

\end{vysvetlenibox}

\newpage

% ============================================================
\section{Akční členy -- elektrické, pneumatické a hydraulické mechanismy}

\begin{pojmybox}
\textbf{Klíčové pojmy:}
servomotor, krokový motor, DC motor, BLDC, solenoid,
přímočarý pohon, pneumatický válec, rozvaděč, hydraulický agregát,
proporcionální ventil, řízení výkonu
\end{pojmybox}

\begin{vysvetlenibox}
\textbf{Elektrické pohony -- srovnání}

\begin{center}
\begin{tabularx}{0.95\linewidth}{@{} l X X X @{}}
    \toprule
    \textbf{Typ motoru} & \textbf{Výhody} & \textbf{Nevýhody}
                        & \textbf{Aplikace} \\
    \midrule
    DC motor       & Jednoduché řízení (napětí = otáčky)
                   & Kartáče -- opotřebení & Modely, ventilátory \\
    BLDC motor     & Bez kartáčů, vysoká účinnost
                   & Složitější řízení (3-fázový) & Drony, e-kola \\
    Krokový motor  & Přesná poloha bez enkodéru
                   & Ztráta kroku při přetížení & 3D tiskárny, CNC \\
    Servomotor AC  & Velmi přesný, rychlý
                   & Drahý, složitý driver & CNC, robotika \\
    \bottomrule
\end{tabularx}
\end{center}

\begin{vzorcebox}
Moment DC motoru: $M = k_t \cdot I_a$,
otáčky: $n = \frac{U_a - I_a R_a}{k_e}$\\
Krokový motor -- krok: $\alpha = \frac{360°}{N_{kroků}}$,
typicky $\alpha = 1{,}8°$ (200 kroků/otáčku)\\
Výkon: $P = M \cdot \omega = M \cdot \frac{2\pi n}{60}$
\end{vzorcebox}

\textbf{Pneumatické systémy}\\
Pracují se stlačeným vzduchem (6--10\,bar).
Výhody: čisté, rychlé, nenákladné.
Nevýhody: hlučné, horší přesnost, energie se špatně zpět získává.

Základní prvky: kompresor $\to$ filtr-regulátor-maznice (FRL)
$\to$ rozvaděč (5/2, 5/3) $\to$ pneumatický válec.

\textbf{Hydraulické systémy}

\begin{vzorcebox}
Pascalův zákon: $p = \frac{F}{A}$ \quad [Pa = N/m²]\\
Síla hydraulického válce: $F = p \cdot A$\\
Příklad: $p = 200\,\text{bar}$, $A = 50\,\text{cm}^2$:
$F = 200 \times 10^5 \cdot 50 \times 10^{-4} = 100\,\text{kN}$
\end{vzorcebox}

\textbf{Srovnání pohonů}

\begin{center}
\begin{tabular}{@{} l c c c @{}}
    \toprule
    \textbf{Vlastnost} & \textbf{Elektrický}
                       & \textbf{Pneumatický} & \textbf{Hydraulický} \\
    \midrule
    Síla / hmotnost & Střední & Nízká & Velmi vysoká \\
    Přesnost polohy & Vysoká & Nízká & Střední \\
    Rychlost        & Vysoká & Velmi vysoká & Střední \\
    Účinnost        & 85--95\,\% & 20--40\,\% & 60--80\,\% \\
    Čistota         & Čistý & Čistý & Riziko úniku oleje \\
    \bottomrule
\end{tabular}

\textbf{Řízení elektrických pohonů}
\begin{itemize}[noitemsep, leftmargin=1.4em, topsep=4pt]
    \item \textbf{PWM regulace} -- změna střídy řídí napětí a tím otáčky
    \item \textbf{H-můstek} -- umožňuje změnu směru otáčení DC motoru
    \item \textbf{FOC (Field Oriented Control)} -- pokročilé řízení BLDC
    \item \textbf{Driver} -- výkonový stupeň mezi MCU a motorem
\end{itemize}

\textbf{Krokový motor -- režimy řízení}
\begin{itemize}[noitemsep, leftmargin=1.4em, topsep=4pt]
    \item Full-step -- celý krok (největší moment)
    \item Half-step -- poloviční krok (vyšší rozlišení)
    \item Microstepping -- plynulý pohyb, menší vibrace
\end{itemize}

\textbf{Servosystém (uzavřená smyčka)}\\
Servomotor pracuje v regulační smyčce:
\begin{itemize}[noitemsep, leftmargin=1.4em, topsep=4pt]
    \item snímač polohy (enkodér)
    \item regulátor (PID)
    \item výkonový měnič
\end{itemize}
Výhoda: vysoká přesnost, žádná ztráta kroku.

\textbf{Pneumatické válce}
\begin{itemize}[noitemsep, leftmargin=1.4em, topsep=4pt]
    \item \textbf{Jednočinný} -- návrat pružinou
    \item \textbf{Dvojčinný} -- tlak na obě strany (větší síla)
    \item \textbf{Tlumení} -- zpomalení na konci zdvihu
\end{itemize}

\textbf{Hydraulika -- hlavní části}
\begin{itemize}[noitemsep, leftmargin=1.4em, topsep=4pt]
    \item čerpadlo -- vytváří tlak
    \item nádrž -- zásoba oleje
    \item ventil -- řízení směru a průtoku
    \item válec / motor -- vykonává práci
\end{itemize}

\textbf{Výhody hydrauliky}
\begin{itemize}[noitemsep, leftmargin=1.4em, topsep=4pt]
    \item extrémně vysoké síly (lisy, bagry)
    \item plynulé řízení
\end{itemize}

\textbf{Nevýhody hydrauliky}
\begin{itemize}[noitemsep, leftmargin=1.4em, topsep=4pt]
    \item únik oleje (ekologie)
    \item složitější údržba
\end{itemize}

\end{center}
\end{vysvetlenibox}

% ============================================================

% ============================================================

\newpage

% ============================================================
\section{Průmyslové roboty a manipulátory}

\begin{pojmybox}
\textbf{Klíčové pojmy:}
DOF, kinematický řetězec, přímá/inverzní kinematika, SCARA,
delta robot, kolaborativní robot (cobot), TCP, trajektorie,
Denavit-Hartenberg, singularita, pracovní prostor
\end{pojmybox}

\begin{vysvetlenibox}

\textbf{Co je průmyslový robot}\\
Průmyslový robot je programovatelný manipulátor, který automaticky
provádí pohybové úlohy (např. svařování, montáž, balení).
Skládá se z mechanické části (rameno), pohonů (motory), senzorů
a řídicí jednotky.

Základní vlastnost robota je \textbf{stupňů volnosti (DOF)}.

\textbf{Stupně volnosti (DOF)}\\
DOF říká, kolik nezávislých pohybů robot má.
Typický průmyslový robot má \textbf{6 DOF}:

\begin{itemize}[noitemsep, leftmargin=1.4em]
    \item 3 DOF pro polohu (X, Y, Z)
    \item 3 DOF pro orientaci (naklopení a otočení nástroje)
\end{itemize}

Díky tomu dokáže robot umístit nástroj do libovolné polohy a směru,
podobně jako lidská ruka.

\textbf{Kinematika robota}\\
Kinematika řeší pohyb robota bez ohledu na síly.

\begin{vzorcebox}

\textbf{Přímá kinematika}\\
Známe úhly kloubů a chceme zjistit polohu nástroje (TCP):

\[
\mathbf{T} = A_1(\theta_1) \cdot A_2(\theta_2) \cdots A_n(\theta_n)
\]

Výstupem je poloha a orientace TCP (Tool Center Point).

\vspace{4pt}

\textbf{Inverzní kinematika}\\
Známe cílovou polohu TCP a hledáme úhly kloubů $\theta_i$.

Tento problém je složitější, protože:
\begin{itemize}[noitemsep]
    \item může existovat více řešení
    \item někdy neexistuje žádné řešení
    \item používají se numerické metody (iterace, Jacobian)
\end{itemize}

\end{vzorcebox}

\textbf{Co je TCP}\\
TCP (Tool Center Point) je bod na konci robota, který se sleduje
a programuje. Například špička svařovacího hořáku nebo uchopovač.

\textbf{Pracovní prostor robota}\\
Je to objem, kam robot dosáhne.
Závisí na:
\begin{itemize}[noitemsep]
    \item délce ramen
    \item konstrukci (SCARA, delta, artikulovaný)
    \item omezeních kloubů
\end{itemize}

\textbf{Typy průmyslových robotů}

\begin{center}
\begin{tabularx}{0.95\linewidth}{@{} l X l @{}}
\toprule
\textbf{Typ} & \textbf{Charakteristika} & \textbf{Použití} \\
\midrule

Artikulovaný (6-DOF) &
Nejvíce podobný lidské ruce, velmi flexibilní,
umožňuje složité pohyby ve 3D prostoru &
Svařování, lakování, manipulace \\

SCARA (4-DOF) &
Rychlý pohyb v rovině XY, velmi tuhý ve vertikále,
ideální pro montáž &
Montáž elektroniky, pick\&place \\

Delta robot (3-DOF) &
Paralelní konstrukce, extrémně rychlý a lehký,
malá hmotnost pohyblivých částí &
Balení, farmaceutický průmysl \\

Kartézský robot (3-DOF) &
Pohyby v osách X, Y, Z (lineární vedení),
jednoduchá konstrukce &
CNC stroje, 3D tisk \\

Cobot (kolaborativní robot) &
Navržen pro práci s člověkem,
má omezení síly a rychlosti pro bezpečnost &
Montáž, asistovaná výroba \\

\bottomrule
\end{tabularx}
\end{center}

\textbf{Trajektorie pohybu}\\
Trajektorie říká, jak se robot pohybuje mezi body:

\begin{itemize}[noitemsep, leftmargin=1.4em]
    \item \textbf{PTP (Point-to-Point)} -- robot jede nejrychlejší cestou,
    dráha mezi body není důležitá
    \item \textbf{CP (Continuous Path)} -- důležitá je přesná dráha
    (např. svařování nebo lakování)
    \item \textbf{Trapézový profil rychlosti} -- pohyb má 3 fáze:
    zrychlení, konstantní rychlost, zpomalení
\end{itemize}

\textbf{Singularita (důležitý pojem)}\\
Singularita nastává, když robot ztrácí jeden stupeň volnosti
a nemůže se pohybovat určitým směrem.
V praxi to může způsobit:
\begin{itemize}[noitemsep]
    \item náhlé zrychlení kloubů
    \item ztrátu kontroly trajektorie
\end{itemize}

\end{vysvetlenibox}

% ============================================================

% ============================================================

\newpage

% ============================================================
\section{Automatizované výrobní a nevýrobní systémy}

\begin{pojmybox}
\textbf{Klíčové pojmy:}
FMS, FMC, AGV, dopravník, zásobník, výrobní buňka,
Industry~4.0, digitální dvojče, MES, ERP, OEE, Lean, Kanban
\end{pojmybox}

\begin{vysvetlenibox}
\textbf{Hierarchie automatizovaných systémů}

\begin{center}
\begin{tabularx}{0.95\linewidth}{@{} l X X @{}}
    \toprule
    \textbf{Systém} & \textbf{Popis} & \textbf{Flexibilita} \\
    \midrule
    Tuhá automatizace & Fixní sled operací (přenosový stroj) & Žádná \\
    FMC (buňka)      & 1--3 CNC stroje + robot + zásobník & Střední \\
    FMS (systém)     & Propojené FMC přes AGV, centrální řízení & Vysoká \\
    \bottomrule
\end{tabularx}
\end{center}

\textbf{Klíčové ukazatele výkonnosti}

\begin{vzorcebox}
OEE (Overall Equipment Effectiveness):
$$OEE = \text{Dostupnost} \times \text{Výkon} \times \text{Kvalita}$$
Příklad: $0{,}90 \times 0{,}85 \times 0{,}99 = 75{,}7\,\%$\\
Světová třída: $OEE > 85\,\%$
\end{vzorcebox}

\textbf{Industry 4.0 -- pilíře}
\begin{itemize}[noitemsep, leftmargin=1.4em, topsep=4pt]
    \item \textbf{IoT} (Internet of Things) -- stroje komunikují přes internet
    \item \textbf{Digitální dvojče} -- virtuální model fyzického zařízení;
          simulace, predikce poruch v reálném čase
    \item \textbf{Cloudová data} -- sběr a analýza dat ve velkém měřítku
    \item \textbf{Kybernetická bezpečnost} -- ochrana řídicích sítí
    \item \textbf{Aditivní výroba} -- 3D tisk součástí na vyžádání
    \item \textbf{Autonomní roboty} -- spolupráce s člověkem, učení
\end{itemize}

\textbf{Struktura výrobní buňky (FMC)}\\
Typická výrobní buňka obsahuje:
\begin{itemize}[noitemsep, leftmargin=1.4em, topsep=4pt]
    \item CNC stroj(e) -- obrábění součástí
    \item Robot -- manipulace s obrobky
    \item Zásobník (magazín) -- ukládání polotovarů a výrobků
    \item Dopravník -- přesun materiálu
    \item PLC -- řízení celé buňky
\end{itemize}

\textbf{Řízení výroby -- MES vs ERP}\\
\begin{itemize}[noitemsep, leftmargin=1.4em, topsep=4pt]
    \item \textbf{ERP} -- plánování výroby, finance, logistika (úroveň firmy)
    \item \textbf{MES} -- řízení výroby v reálném čase (zakázky, stroje, data)
    \item \textbf{SCADA} -- vizualizace a dohled nad procesem
\end{itemize}

\textbf{Lean manufacturing (štíhlá výroba)}\\
Cílem je eliminace plýtvání:
\begin{itemize}[noitemsep, leftmargin=1.4em, topsep=4pt]
    \item nadprodukce
    \item čekání
    \item zbytečný pohyb
    \item vady (zmetky)
    \item nadbytečné zásoby
\end{itemize}

\textbf{Výhody automatizace}
\begin{itemize}[noitemsep, leftmargin=1.4em, topsep=4pt]
    \item vyšší produktivita a kvalita
    \item menší chybovost (eliminace lidského faktoru)
    \item možnost nepřetržitého provozu (24/7)
\end{itemize}

\textbf{Nevýhody automatizace}
\begin{itemize}[noitemsep, leftmargin=1.4em, topsep=4pt]
    \item vysoké pořizovací náklady
    \item složitější údržba a servis
    \item menší flexibilita u tuhých systémů
\end{itemize}

\textbf{AGV} (Automated Guided Vehicle) -- autonomní vozidla pro dopravu
materiálu; navigace: laserová, magnetická dráha, vizuální (kamera + AI).
\end{vysvetlenibox}

% ============================================================

% ============================================================


\newpage


% ============================================================
\section{Umělá inteligence}

\begin{pojmybox}
strojové učení, neuronová síť, hluboké učení, CNN, RNN, transformer,
predikátová logika, fuzzy logika, inference, dataset, trénování modelu
\end{pojmybox}

% ============================================================
\begin{vysvetlenibox}

\textbf{Základní rozdělení AI}

\begin{itemize}
\item \textbf{Symbolická AI} – expertní systémy, pravidla IF-THEN
\item \textbf{Datová AI} – strojové učení, neuronové sítě
\end{itemize}

% ============================================================
\textbf{Neuron – základní stavební jednotka AI}

Neuron simuluje chování biologického neuronu. Zpracovává vstupy,
váží je, přičítá bias a aplikuje aktivační funkci.

\begin{center}
\begin{tikzpicture}[node distance=1.3cm]

\node (x1) at (0,1.2) {$x_1$};
\node (x2) at (0,0.4) {$x_2$};
\node (x3) at (0,-0.4) {$x_3$};

\node[circle, draw, minimum size=1.4cm] (sum) at (2,0.4) {$\Sigma$};

\node (f) at (3.8,0.4) {f(x)};
\node (y) at (5.2,0.4) {$y$};

\draw[->] (x1) -- node[above] {$w_1$} (sum);
\draw[->] (x2) -- node[above] {$w_2$} (sum);
\draw[->] (x3) -- node[below] {$w_3$} (sum);

\draw[->] (sum) -- node[above] {$+b$} (f);
\draw[->] (f) -- (y);

\end{tikzpicture}
\end{center}

\begin{vzorcebox}
\[
y = f\left(\sum_{i=1}^{n} w_i x_i + b\right)
\]
\end{vzorcebox}

\textbf{Co neuron obsahuje:}
\begin{itemize}
\item vstupy $x_i$ – data (např. senzory)
\item váhy $w_i$ – důležitost vstupu
\item bias $b$ – posun rozhodnutí
\item sumační funkce $\Sigma$
\item aktivační funkce $f(x)$
\end{itemize}

% ============================================================
\textbf{Aktivační funkce}

\begin{itemize}
\item ReLU: $f(x)=\max(0,x)$
\item Sigmoid: $f(x)=\frac{1}{1+e^{-x}}$
\item Tanh: škála -1 až 1
\end{itemize}

% ============================================================
\textbf{Neuronová síť}

Více neuronů tvoří vrstvy:
\begin{itemize}
\item vstupní vrstva
\item skryté vrstvy
\item výstupní vrstva
\end{itemize}

Učí se pomocí zpětné propagace chyby (backpropagation).

% ============================================================
\textbf{Predikátová logika}

Pracuje s proměnnými a kvantifikátory:

\[
\forall x \, (\text{člověk}(x) \Rightarrow \text{smrtelný}(x))
\]

Použití:
\begin{itemize}
\item expertní systémy
\item automatické dokazování
\end{itemize}

% ============================================================
\textbf{Fuzzy logika}

Na rozdíl od klasické logiky nepracuje jen s 0/1.

\[
\mu(x) \in [0,1]
\]

Příklad:
\begin{itemize}
\item teplota = "vysoká" → 0.7
\item rychlost = "střední" → 0.5
\end{itemize}

Použití:
\begin{itemize}
\item řízení klimatizace
\item pračky, robotika
\end{itemize}

% ============================================================
\textbf{ML vs expertní systémy}

\begin{tabularx}{\linewidth}{l X X}
\toprule
 & Strojové učení & Expertní systémy \\
\midrule
Znalosti & učí se z dat & ručně zadaná pravidla \\
Učení & ano & ne \\
Flexibilita & vysoká & nízká \\
Úprava & automatická & manuální \\
Použití & predikce, klasifikace & rozhodovací pravidla \\
\bottomrule
\end{tabularx}

% ============================================================
\end{vysvetlenibox}



\newpage


\section{Nové materiály a technologie v mechatronice}

\begin{pojmybox}
\textbf{Klíčové pojmy:}
kompozity, uhlíková vlákna, SMA, piezoelektrika,
MEMS, nanomateriály, 3D tisk,
FDM, SLA, SLS, DMLS, smart materiály
\end{pojmybox}

\begin{vysvetlenibox}

% ------------------------------------------------------------
\textbf{Konstrukční materiály -- srovnání}

\begin{center}
\begin{tabular}{l r r r r}
Materiál & $\rho$ & $E$ & $R_m$ & $E/\rho$ \\
 & [g/cm³] & [GPa] & [MPa] & \\
\hline
Ocel S235 & 7,85 & 210 & 360 & 26,8 \\
Al 7075   & 2,81 & 72  & 570 & 25,6 \\
CFRP      & 1,55 & 135 & 1500 & 87,1 \\
Titan     & 4,43 & 114 & 950 & 25,7 \\
\end{tabular}
\end{center}

\textbf{Shrnutí:}
\begin{itemize}
\item Kompozity (CFRP) mají nejlepší poměr pevnost/hmotnost
\item Ocel je levná a pevná, ale těžká
\item Hliník je lehký, ale méně tuhý
\end{itemize}

% ------------------------------------------------------------
\textbf{Smart materiály}

\begin{itemize}
\item \textbf{SMA (tvarová paměť)}
    \begin{itemize}
    \item materiál si „pamatuje“ původní tvar
    \item při zahřátí se vrací zpět
    \item použití: aktuátory, medicína
    \end{itemize}

\item \textbf{Piezoelektrika}
    \begin{itemize}
    \item tlak → napětí (senzor)
    \item napětí → deformace (aktuátor)
    \item velmi rychlá odezva
    \end{itemize}

\begin{vzorcebox}
Piezo jev: $Q = d \cdot F$
\end{vzorcebox}

\item \textbf{Magnetoreologické kapaliny}
    \begin{itemize}
    \item mění viskozitu v magnetickém poli
    \item použití: tlumiče (auta, robotika)
    \end{itemize}
\end{itemize}

% ------------------------------------------------------------
\textbf{MEMS (Micro-Electro-Mechanical Systems)}

Miniaturní mechanické struktury na čipu:

\begin{itemize}
\item akcelerometry (mobil, drony)
\item gyroskopy
\item tlakové senzory
\end{itemize}

\textbf{Výhody:}
\begin{itemize}
\item malé rozměry
\item nízká cena
\item integrace s elektronikou
\end{itemize}

% ------------------------------------------------------------
\textbf{Nanomateriály}

\begin{itemize}
\item struktury v měřítku nm ($10^{-9}$ m)
\item extrémní vlastnosti:
\begin{itemize}
\item vysoká pevnost
\item vysoká vodivost
\end{itemize}
\item příklad: grafen, nanotrubice
\end{itemize}

% ------------------------------------------------------------
\textbf{3D tisk (aditivní výroba)}

\begin{center}
\begin{tabularx}{0.95\linewidth}{l X l}
Metoda & Princip & Materiál \\
\hline
FDM  & tavení plastu & PLA, ABS \\
SLA  & UV vytvrzení & pryskyřice \\
SLS  & spékání prášku & nylon \\
DMLS & tavení kovu & ocel, titan \\
\end{tabularx}
\end{center}

\textbf{Výhody:}
\begin{itemize}
\item složité tvary bez obrábění
\item rychlé prototypování
\end{itemize}

\textbf{Nevýhody:}
\begin{itemize}
\item horší mechanické vlastnosti (hlavně FDM)
\item vyšší cena u kovového tisku
\end{itemize}

% ------------------------------------------------------------
\textbf{Shrnutí}

\begin{itemize}
\item Kompozity → lehké a pevné konstrukce
\item Smart materiály → senzory + aktuátory
\item MEMS → miniaturizace
\item 3D tisk → rychlá výroba
\end{itemize}

\end{vysvetlenibox}



\newpage


\section{Metody technické diagnostiky}

\begin{pojmybox}
\textbf{Klíčové pojmy:}
prediktivní údržba, vibrodiagnostika, termodiagnostika,
tribodiagnostika, akustická emise, NDT,
spolehlivost, MTBF, MTTR, Weibullovo rozdělení
\end{pojmybox}

\begin{vysvetlenibox}

% ------------------------------------------------------------
\textbf{Strategie údržby}

\begin{center}
\begin{tabularx}{0.95\linewidth}{l X X}
Strategie & Popis & Náklady \\
\hline
Reaktivní & oprava až po poruše & velmi vysoké \\
Preventivní & pravidelná výměna & střední \\
Prediktivní & dle stavu zařízení & nejnižší \\
\end{tabularx}
\end{center}

\textbf{Shrnutí:}
\begin{itemize}
\item cílem je přejít z reaktivní → prediktivní údržbu
\item využívají se senzory + analýza dat
\end{itemize}

% ------------------------------------------------------------
\textbf{Ukazatele spolehlivosti}

\begin{vzorcebox}
$MTBF = \frac{t_{provoz}}{N_{poruch}}$

$MTTR = \frac{t_{opravy}}{N_{oprav}}$

$A = \frac{MTBF}{MTBF + MTTR}$
\end{vzorcebox}

\textbf{Význam:}
\begin{itemize}
\item MTBF = jak často se to rozbije
\item MTTR = jak dlouho to opravuješ
\item dostupnost = jak moc je stroj „k dispozici“
\end{itemize}

% ------------------------------------------------------------
\textbf{Vibrodiagnostika}

\begin{itemize}
\item měření vibrací (akcelerometr)
\item analýza frekvence (FFT)
\item každá porucha má typický signál
\end{itemize}

\begin{center}
\begin{tabularx}{0.95\linewidth}{l X}
Závada & Projev \\
\hline
Nevyváženost & 1× otáčky \\
Nesouosost & 2× otáčky \\
Ložiska & specifické frekvence (BPFO, BPFI) \\
Ozubená kola & násobky zubů \\
\end{tabularx}
\end{center}

% ------------------------------------------------------------
\textbf{Termodiagnostika}

\begin{itemize}
\item měření teploty (IR kamera)
\item odhaluje:
\begin{itemize}
\item přehřáté ložiska
\item špatné kontakty (elektrika)
\item tření
\end{itemize}
\end{itemize}

% ------------------------------------------------------------
\textbf{Tribodiagnostika}

Analýza oleje:

\begin{itemize}
\item obsah kovových částic → opotřebení
\item viskozita → degradace oleje
\item voda → porucha těsnění
\end{itemize}

% ------------------------------------------------------------
\textbf{NDT metody (nedestruktivní testování)}

\begin{itemize}
\item RT (rentgen) → vnitřní vady
\item UT (ultrazvuk) → trhliny, tloušťka
\item MT → feromagnetické materiály
\item PT → povrchové trhliny
\end{itemize}

% ------------------------------------------------------------
\textbf{Weibullovo rozdělení}

Používá se pro modelování poruch:

\begin{itemize}
\item rané poruchy (výrobní vady)
\item náhodné poruchy
\item opotřebení (stárnutí)
\end{itemize}

\textbf{Tzv. vanová křivka (bathtub curve):}
\begin{itemize}
\item začátek → vysoká poruchovost
\item střed → stabilní provoz
\item konec → opotřebení
\end{itemize}

% ------------------------------------------------------------
\textbf{Shrnutí}

\begin{itemize}
\item cílem diagnostiky je zabránit poruchám
\item využívá se měření (vibrace, teplota, olej)
\item moderní trend = prediktivní údržba + data
\end{itemize}

\end{vysvetlenibox}

\newpage

% ============================================================
\section{Převody točivého pohybu (ozubené, řetězové, řemenové)}

\begin{pojmybox}
\textbf{Klíčové pojmy:}
převodový poměr, momentový poměr, účinnost, ozubené kolo,
modul, řemen (plochý, klínový, ozubený), řetěz,
vůle, mazání, evolventní profil, čelní/kuželové ozubení
\end{pojmybox}

\begin{vysvetlenibox}
\textbf{Základní vztahy}

\begin{vzorcebox}
Převodový poměr: $i = \frac{n_1}{n_2} = \frac{z_2}{z_1} = \frac{d_2}{d_1}$\\
Momentový přenos: $M_2 = M_1 \cdot i \cdot \eta$
\quad ($\eta$ = účinnost soukolí)\\
Průměr roztečné kružnice: $d = m \cdot z$
\quad (modul $m$ [mm], počet zubů $z$)\\
Osová vzdálenost čelního soukolí:
$a = \frac{m(z_1 + z_2)}{2}$
\end{vzorcebox}

\textbf{Typy ozubených kol}

\begin{center}
\begin{tabularx}{0.95\linewidth}{@{} l X l @{}}
    \toprule
    \textbf{Typ} & \textbf{Vlastnosti} & \textbf{Účinnost} \\
    \midrule
    Čelní přímé    & Jednoduché, hlučnější, riziko kavitace & 97--99\,\% \\
    Čelní šikmé    & Tišší, větší axiální síla & 97--99\,\% \\
    Kuželové       & Změna osy rotace (90°) & 96--98\,\% \\
    Šneková        & Velký převodový poměr, samosvornost & 50--90\,\% \\
    Planetové      & Kompaktní, velký moment, koaxiální & 95--98\,\% \\
    \bottomrule
\end{tabularx}
\end{center}

\textbf{Řemenové a řetězové převody}

\begin{center}
\begin{tabularx}{0.95\linewidth}{@{} l X X X @{}}
    \toprule
    & \textbf{Klínový řemen} & \textbf{Ozubený řemen}
    & \textbf{Řetěz} \\
    \midrule
    Skluz      & Ano (3--5\,\%) & Ne & Ne \\
    Tlumení    & Vysoké & Střední & Nízké \\
    Mazání     & Ne & Ne & Ano \\
    Účinnost   & 92--96\,\% & 96--99\,\% & 96--98\,\% \\
    Použití    & Kompresory, ventilátory & Rozvod, servopohony
               & Motorky, pohony os \\
    \bottomrule
\end{tabularx}
\end{center}

\textbf{Základní typy převodovek}
\begin{center}
\begin{tabular}{|l|c|c|c|c|}
\hline
 & Běžná (čelní) & Planetární & Harmonická & Cykloidní \\
\hline
Princip & Ozubená kola & Kolo + satelity & Pružná deformace & Excentrický kotouč \\
\hline
Převod & Nízký--střední & Střední--vysoký & Velmi vysoký & Velmi vysoký \\
\hline
Kompaktnost & Nízká & Vysoká & Velmi vysoká & Vysoká \\
\hline
Vůle & Střední & Nízká & Velmi nízká & Velmi nízká \\
\hline
Účinnost & 95--98\,\% & 95--97\,\% & 75--90\,\% & 85--95\,\% \\
\hline
Použití & Stroje & Robotika & Přesné pohony & Reduktory \\
\hline
\end{tabular}
\end{center}

\end{vysvetlenibox}

% ============================================================

% ============================================================

\newpage

% ============================================================
\section{Mechanismy obecného pohybu a tvrdá automatizace}

\begin{pojmybox}
\textbf{Klíčové pojmy:}
vačka, páka, kliková hřídel, kulisový mechanismus, Genevský mechanismus,
tvrdá automatizace, přenosový stroj, karuselový stroj, stupeň volnosti,
Grüblerova podmínka
\end{pojmybox}

\begin{vysvetlenibox}
\textbf{Stupně volnosti mechanismu}

\begin{vzorcebox}
Grüblerova podmínka (rovinný mechanismus):
$$F = 3(n-1) - 2p_1 - p_2$$
kde $n$ = počet členů (včetně rámu), $p_1$ = počet jednoduchých kloubů,
$p_2$ = počet dvojkloubů. Mechanismus: $F=1$, struktura: $F=0$.
\end{vzorcebox}

\textbf{Základní mechanismy}

\begin{center}
\begin{tabularx}{0.95\linewidth}{@{} l X l @{}}
    \toprule
    \textbf{Mechanismus} & \textbf{Funkce} & \textbf{Aplikace} \\
    \midrule
    Vačkový      & Rotace $\to$ libovolný pohyb sledovače;
                   profil vačky určuje zákon pohybu & Ventily motoru, podavače \\
    Klikový       & Rotace $\to$ přímočarý pohyb & Pístové motory, lisy \\
    Genevský      & Přerušovaný pohyb (1 otočení $\to$ 1 krok)
                  & Indexovací stoly, filmové projektory \\
    Páka          & Přenos a změna síly/pohybu & Záchyty, spínače \\
    Kulisový      & Kluzná i rotační vazba & Variátory rychlosti \\
    \bottomrule
\end{tabularx}
\end{center}

\textbf{Tvrdá automatizace -- charakteristika}\\
Stroje pro hromadnou výrobu s neměnným sledem operací.
Maximální produktivita, ale nulová flexibilita při změně výrobku.

\begin{center}
\begin{tabularx}{0.95\linewidth}{@{} l X @{}}
    \toprule
    \textbf{Typ stroje} & \textbf{Popis} \\
    \midrule
    Přenosový (in-line) & Obrobky se posunují lineárně od stanice ke stanici \\
    Karuselový (rotační) & Otočný stůl, obrobky se otáčejí kolem osy \\
    Gnezdový (cellular)  & Obrobky se zpracovávají na jednom místě \\
    \bottomrule
\end{tabularx}
\end{center}
\end{vysvetlenibox}

% ============================================================

% ============================================================

\newpage

% ============================================================
\section{Číslicově řízené stroje (CNC) a jejich programování}

\begin{pojmybox}
\textbf{Klíčové pojmy:}
G-kód, M-kód, souřadnicový systém, absolutní/přírůstkové programování,
nulový bod, korekce nástroje, frézování, soustružení, CAM,
interpolace, posuvová rychlost, otáčky vřetena
\end{pojmybox}

\begin{vysvetlenibox}
\textbf{Princip CNC stroje}\\
Existují primárně 3 druhy CNC strojů - 3 osy, 4 osy, 5 os, a jejich kombinace (např 3+1 znamený tří-osý se čtvrtou osou přidanou)
\textbf{Souřadnicový systém CNC stroje}\\
Pravotočivý kartézský systém (X, Y, Z + rotační osy A, B, C).
Nulový bod obrobku (WPZ) se nastavuje před obráběním. Nulový bod stroje (MZ) se nastaví hned po zapnutí stroje například koncovými spínači a bez jeho nastavení jede stroj na slepo, jelikož neví kde se nachází.

\textbf{Přehled G-kódů}

\begin{center}
\begin{tabularx}{0.95\linewidth}{@{} l X @{}}
    \toprule
    \textbf{G-kód} & \textbf{Funkce} \\
    \midrule
    G00 & Rychloposuv (bez řezu) \\
    G01 & Lineární interpolace (řez přímou čarou) \\
    G02 & Kruhová interpolace po směru hodinových ručiček \\
    G03 & Kruhová interpolace proti směru hod. ručiček \\
    G17/18/19 & Výběr roviny (XY / XZ / YZ) \\
    G90/G91 & Absolutní / přírůstkové souřadnice \\
    G94/G95 & Posuv [mm/min] / [mm/otáčku] \\
    \bottomrule
\end{tabularx}
\end{center}

\textbf{Přehled M-kódů}

\begin{center}
\begin{tabular}{@{} l l @{\quad} l l @{}}
    \toprule
    M00 & Pauza & M08 & Chlazení zapnout \\
    M03 & Vřeteno CW & M09 & Chlazení vypnout \\
    M04 & Vřeteno CCW & M30 & Konec programu \\
    M05 & Vřeteno stop & M06 & Výměna nástroje \\
    \bottomrule
\end{tabular}
\end{center}

\textbf{Ukázkový blok G-kódu:}
\begin{vzorcebox}
\begin{verbatim}
N10 G90 G17 G94         ; Absolutní, rovina XY, mm/min
N20 T01 M06             ; Nástroj č. 1
N30 S1200 M03           ; 1200 RPM, vřeteno CW
N40 G00 X0 Y0 Z5        ; Rychloposuv nad nulový bod
N50 G01 Z-3 F100        ; Hloubkový záběr 3mm, F=100
N60 G01 X50 Y0 F300     ; Frézování přímkou X+50
N70 G02 X50 Y50 I0 J25  ; Oblouk CW, střed I0 J25
N80 G00 Z50             ; Zdvih nástroje
N90 M05 M30             ; Stop vřetena, konec
\end{verbatim}
\end{vzorcebox}

\textbf{CAM -- Computer Aided Manufacturing}\\
CAM software (Fusion 360, Mastercam, NX, SolidCAM, FreeCAD CAM) generuje G-kód z nakreslených/navrhlých cest. Postprocesor převede neutrální CLData na G-kód specifický
pro daný řídicí systém (Fanuc, Siemens Sinumerik, Heidenhain, GRBL, Hase).

\begin{vzorcebox}
Řezná rychlost: $v_c = \frac{\pi \cdot d \cdot n}{1000}$ [m/min]\\
$\Rightarrow$ Otáčky: $n = \frac{1000 \cdot v_c}{\pi \cdot d}$ [min$^{-1}$]\\
Minutový posuv: $v_f = f_z \cdot z \cdot n$ [mm/min]
\quad ($f_z$ = posuv na zub, $z$ = počet zubů)
\end{vzorcebox}
\end{vysvetlenibox}

% ============================================================

% ============================================================

\newpage

% ============================================================
\section{CAx systémy a Rapid prototyping}

\begin{pojmybox}
\textbf{Klíčové pojmy:}
CAD, CAM, CAE, FEM simulace, Fusion~360, SolidWorks, CATIA,
Rapid prototyping, FDM, SLA, SLS, DMLS, postprocesor,
STL formát, 3MF formát, optimalizace topologie
\end{pojmybox}

\begin{vysvetlenibox}
\textbf{CAx -- přehled nástrojů}

\begin{center}
\begin{tabularx}{0.95\linewidth}{@{} l l X @{}}
    \toprule
    \textbf{Zkratka} & \textbf{Název} & \textbf{Funkce} \\
    \midrule
    CAD & Computer Aided Design      & 3D parametrické modelování geometrie \\
    CAM & Computer Aided Manufacturing & Návrh a generování gcode drah z modelu za použití postprocesoru\\
    CAE & Computer Aided Engineering & Simulace a analýza (FEM, CFD, MBS) \\
    FEM & Finite Element Method      & Simulace pevnosti, slabých bodů a ohybu modelu \\
    CAPP& Aided Process Planning     & Plánování výrobního procesu \\
    PDM & Product Data Management    & Správa dat a verzí modelu \\
    \bottomrule
\end{tabularx}
\end{center}

\textbf{Metoda konečných prvků (FEM)}\\
Z modelu se vytvoří mřížka z mnoha trojúhelníků (mesh). Pro každý trojúhelník se zapíše
rovnice tuhosti. Soustava rovnic se řeší numericky. Přidají se tažné, tlačné, střihové a otočné síly, a hledají se ohniska tažných a tlačivých sil či vysokou lokální zátěží (neboli slabé body).

\begin{vzorcebox}
Maticová forma: $[K]\{u\} = \{F\}$\\
kde $[K]$ = globální matice tuhosti, $\{u\}$ = posuvy uzlů,
$\{F\}$ = vektor sil.\\
Napětí: $\sigma = E \cdot \varepsilon = E \cdot \frac{\Delta l}{l_0}$
\end{vzorcebox}

\textbf{Rapid prototyping -- pracovní postup}

\begin{center}
\begin{tikzpicture}[scale=0.85,
  step/.style={rectangle, draw=accent, fill=lightblue, rounded corners=3pt,
               minimum width=2.3cm, minimum height=0.6cm,
               text centered, font=\small}]
  \node[inner sep=2pt,rectangle,draw,fill=blue!20] (cad)  at (-2,0)    {3D model (CAD) a FEM};
  \node[inner sep=2pt,rectangle,draw,fill=blue!20] (stl)  at (3,0)    {Export STL/3MF};
  \node[inner sep=2pt,rectangle,draw,fill=blue!20] (slic) at (6,0)    {Slicování};
  \node[inner sep=2pt,rectangle,draw,fill=blue!20] (tisk) at (6,-1.2) {Tisk / výroba};
  \node[inner sep=2pt,rectangle,draw,fill=blue!20] (post) at (3,-1.2) {Post-processing};
  \node[inner sep=2pt,rectangle,draw,fill=blue!20] (test) at (0,-1.2) {Testování};
  \draw[>-,accent] (cad)--(stl)--(slic)--(tisk)--(post)--(test);
\end{tikzpicture}
\end{center}

\textbf{Topologická optimalizace}\\
Algoritmus odstraní materiál z míst s nízkým napětím -- výsledek
je organická struktura s maximálním poměrem tuhosti/hmotnosti.
Typicky se kombinuje s DMLS tiskem (nelze vyrobit klasickým obráběním).

\textbf{Srovnání RP technologií -- použití v praxi}

\begin{center}
\begin{tabularx}{0.95\linewidth}{@{} l X X @{}}
    \toprule
    \textbf{Metoda} & \textbf{Výhody} & \textbf{Nevýhody} \\
    \midrule
    FDM  & Levná, dostupná, velký sortiment materiálů
         & Viditelné vrstvy, nižší přesnost, každá vrstva je potencionální slabý bod \\
    SLA  & Hladký povrch, vysoká přesnost, potencionálně rychlá při paralelním tisku
         & Drahá pryskyřice, křehké \\
    SLS  & Bez podpor, funkční díly z nylonu
         & Drahý stroj, prašné \\
    DMLS & Kovové díly, složité tvary, vysoká pevnost
         & Velmi drahý stroj (\textgreater{}500k EUR) \\
    \bottomrule
\end{tabularx}
\end{center}
\end{vysvetlenibox}

% ============================================================
\vfill
\noindent\textcolor{medgray}{\small
Souhlasím se zpracováním osobních údajů pro účely náboru v souladu s GDPR.}


% ============================================================
\end{document}
